% Options for packages loaded elsewhere
\PassOptionsToPackage{unicode}{hyperref}
\PassOptionsToPackage{hyphens}{url}
\documentclass[
]{article}
\usepackage{xcolor}
\usepackage[margin=1in]{geometry}
\usepackage{amsmath,amssymb}
\setcounter{secnumdepth}{-\maxdimen} % remove section numbering
\usepackage{iftex}
\ifPDFTeX
  \usepackage[T1]{fontenc}
  \usepackage[utf8]{inputenc}
  \usepackage{textcomp} % provide euro and other symbols
\else % if luatex or xetex
  \usepackage{unicode-math} % this also loads fontspec
  \defaultfontfeatures{Scale=MatchLowercase}
  \defaultfontfeatures[\rmfamily]{Ligatures=TeX,Scale=1}
\fi
\usepackage{lmodern}
\ifPDFTeX\else
  % xetex/luatex font selection
\fi
% Use upquote if available, for straight quotes in verbatim environments
\IfFileExists{upquote.sty}{\usepackage{upquote}}{}
\IfFileExists{microtype.sty}{% use microtype if available
  \usepackage[]{microtype}
  \UseMicrotypeSet[protrusion]{basicmath} % disable protrusion for tt fonts
}{}
\makeatletter
\@ifundefined{KOMAClassName}{% if non-KOMA class
  \IfFileExists{parskip.sty}{%
    \usepackage{parskip}
  }{% else
    \setlength{\parindent}{0pt}
    \setlength{\parskip}{6pt plus 2pt minus 1pt}}
}{% if KOMA class
  \KOMAoptions{parskip=half}}
\makeatother
\usepackage{longtable,booktabs,array}
\newcounter{none} % for unnumbered tables
\usepackage{calc} % for calculating minipage widths
% Correct order of tables after \paragraph or \subparagraph
\usepackage{etoolbox}
\makeatletter
\patchcmd\longtable{\par}{\if@noskipsec\mbox{}\fi\par}{}{}
\makeatother
% Allow footnotes in longtable head/foot
\IfFileExists{footnotehyper.sty}{\usepackage{footnotehyper}}{\usepackage{footnote}}
\makesavenoteenv{longtable}
\usepackage{graphicx}
\makeatletter
\newsavebox\pandoc@box
\newcommand*\pandocbounded[1]{% scales image to fit in text height/width
  \sbox\pandoc@box{#1}%
  \Gscale@div\@tempa{\textheight}{\dimexpr\ht\pandoc@box+\dp\pandoc@box\relax}%
  \Gscale@div\@tempb{\linewidth}{\wd\pandoc@box}%
  \ifdim\@tempb\p@<\@tempa\p@\let\@tempa\@tempb\fi% select the smaller of both
  \ifdim\@tempa\p@<\p@\scalebox{\@tempa}{\usebox\pandoc@box}%
  \else\usebox{\pandoc@box}%
  \fi%
}
% Set default figure placement to htbp
\def\fps@figure{htbp}
\makeatother
\setlength{\emergencystretch}{3em} % prevent overfull lines
\providecommand{\tightlist}{%
  \setlength{\itemsep}{0pt}\setlength{\parskip}{0pt}}
\usepackage{graphicx}
\usepackage{array}
\makeatletter
\renewcommand{\maketitle}{%
  \begin{titlepage}
  \thispagestyle{empty}
  \centering
  \vspace*{1.3cm}
  {\bfseries\fontsize{16}{20}\selectfont
  AUDIT REPRODUCIBILITY DAN KUALITAS DATA\\
  PADA REGRESI BAYESIAN UNTUK PREDIKSI\\
  HAMBATAN RESIDU KAPAL LAYAR\\
  (STUDI KASUS: YACHT HYDRODYNAMICS)\par}
  \vspace{1.0cm}
  {\fontsize{12}{16}\selectfont\itshape
  Ditujukan untuk Memenuhi Tugas Project Mata Kuliah Analisis Regresi Tingkat Lanjut\par}
  \vspace{0.9cm}
  {\fontsize{12}{16}\selectfont Dosen Pengampu:\par}
  \vspace{0.2cm}
  {\bfseries\fontsize{12}{16}\selectfont Dr. I Gede Nyoman Mindra Jaya, S.Si., M.Si.\par}
  \vspace{0.8cm}
  \IfFileExists{logo-unpad1.png}{%
    \includegraphics[width=0.22\textwidth]{logo-unpad1.png}\par
  }{% 
    {\bfseries\fontsize{12}{16}\selectfont UNIVERSITAS PADJADJARAN\par}
  }
  \vspace{0.8cm}
  {\fontsize{12}{16}\selectfont Disusun Oleh:\par}
  \vspace{0.2cm}
  {\bfseries\fontsize{12}{16}\selectfont
  Yaafi Ferdian Syahputra (140720260018)\\
  I Gusti Bagus Harya Putra (140720260016)\par}
  \vspace{0.9cm}
  {\bfseries\fontsize{12}{16}\selectfont
  PROGRAM STUDI MAGISTER STATISTIKA TERAPAN\\
  DEPARTEMEN STATISTIKA\\
  FAKULTAS MATEMATIKA DAN ILMU PENGETAHUAN ALAM\\
  UNIVERSITAS PADJADJARAN\\
  TAHUN 2026\par}
  \vfill
  \end{titlepage}
}
\renewcommand{\contentsname}{DAFTAR ISI}
\makeatother
\usepackage{booktabs}
\usepackage{longtable}
\usepackage{array}
\usepackage{multirow}
\usepackage{wrapfig}
\usepackage{float}
\usepackage{colortbl}
\usepackage{pdflscape}
\usepackage{tabu}
\usepackage{threeparttable}
\usepackage{threeparttablex}
\usepackage[normalem]{ulem}
\usepackage{makecell}
\usepackage{xcolor}
\usepackage{bookmark}
\IfFileExists{xurl.sty}{\usepackage{xurl}}{} % add URL line breaks if available
\urlstyle{same}
% fallback for those not using the hyperref driver hyperxmp:
\makeatletter
\@ifundefined{xmpquote}{\newcommand{\xmpquote}[1]{#1}}{}
\makeatother
\hypersetup{
  pdftitle={AUDIT REPRODUCIBILITY DAN KUALITAS DATA PADA REGRESI BAYESIAN UNTUK PREDIKSI HAMBATAN RESIDU KAPAL LAYAR},
  pdfauthor={I Gusti Bagus Harya Putra dan Yaafi Ferdian Syahputra},
  hidelinks,
  pdfcreator={LaTeX via pandoc}}

\title{AUDIT REPRODUCIBILITY DAN KUALITAS DATA PADA REGRESI BAYESIAN
UNTUK PREDIKSI HAMBATAN RESIDU KAPAL LAYAR}
\author{I Gusti Bagus Harya Putra dan Yaafi Ferdian Syahputra}
\date{2026}

\begin{document}
\maketitle

\renewcommand*\contentsname{DAFTAR ISI}
{
\setcounter{tocdepth}{2}
\tableofcontents
}
\textbf{ABSTRAK}

Prediksi hambatan residu (\emph{residuary resistance}) kapal layar pada
tahap desain awal penting untuk menaksir kebutuhan daya propulsi.
Penelitian ini menjawab pertanyaan: \emph{seberapa stabil dan informatif
hasil Regresi Bayesian pada data Yacht Hydrodynamics ketika difokuskan
pada reproducibility dan kualitas data?} Data terdiri atas 308 uji
\emph{towing tank} dari 22 bentuk lambung × 14 Froude number. Analisis
mengikuti \emph{Bayesian workflow}: audit data, prior predictive check,
validasi komputasi (analitik, Gibbs, dan HMC), tujuh model bertahap
(M0--M6), diagnostik MCMC, \emph{posterior predictive check}, PSIS-LOO
dengan koreksi Jacobian, analisis sensitivitas, serta audit
reproducibility (seed, urutan data, dan skema validasi). Komputasi
\textbf{stabil secara substantif}: posterior MCMC cocok dengan solusi
analitik, dan probabilitas eksedans prediktif hanya bervariasi
0.757--0.765 antar lima seed. Uji formal \textbar z\textbar{}
\textless{} 3 MCSE menandai SD efek hull; audit lima seed menunjukkan
penyebabnya adalah MCSE yang meremehkan variasi antar-run (≈0.5×), bukan
non-determinisme. Sebaliknya, keputusan \textbf{kualitas data dan
preprocessing mendominasi hasil}: model Gaussian pada skala mentah (M0)
konvergen sempurna namun salah spesifikasi (27\% replikasi bernilai
negatif), dan RMSE validasi per-hull turun dari 8.87 (M0) menjadi 1.87
(M3). Validasi acak per baris \textbf{melebih-lebihkan} kinerja model
hierarkis sebesar 63.4 ELPD (SE 11.7) dan mengubah model yang terpilih.
Model informatif untuk prediksi (cakupan interval prediksi 95\% per-hull
M3 = 0.94; M6 sedikit konservatif, 0.99), tetapi arah efek masing-masing
rasio geometri tidak dapat dipastikan dari 22 hull.

\textbf{Kata Kunci:} Regresi Bayesian, Reproducibility, Kualitas Data,
PSIS-LOO, \emph{Data Leakage}, Yacht Hydrodynamics.

\clearpage

\section{BAB 1. PENDAHULUAN}\label{bab-1.-pendahuluan}

\subsection{1.1 Latar Belakang Masalah}\label{latar-belakang-masalah}

Hambatan residu per satuan berat displacement adalah komponen hambatan
kapal di luar gesekan, terutama akibat pembentukan gelombang. Nilainya
menentukan kebutuhan daya propulsi, sehingga prediksi yang akurat pada
tahap desain awal sangat bernilai. Dataset \emph{Delft Systematic Yacht
Hull Series} (UCI ID 243) memuat hasil uji \emph{towing tank} untuk 22
bentuk lambung yang diturunkan dari satu \emph{parent hull} (Standfast
43), masing-masing diuji pada 14 kecepatan (Froude number).

Regresi Bayesian menawarkan kuantifikasi ketidakpastian penuh melalui
distribusi posterior dan posterior predictive. Namun, hasil Bayesian
\textbf{bersyarat pada model, prior, data, dan komputasi}. Kesalahan
kecil pada kualitas data atau preprocessing --- misalnya skala prior
yang tidak sesuai, transformasi target, atau skema validasi yang bocor
--- dapat menghasilkan kesimpulan yang tampak meyakinkan (R̂ = 1.00,
tanpa divergensi) tetapi keliru. Karena itu, fokus project ini adalah
\textbf{reproducibility dan kualitas data}: seberapa stabil hasil
terhadap sumber variasi yang seharusnya tidak berpengaruh, dan seberapa
besar pengaruh keputusan data terhadap kesimpulan.

\subsection{1.2 Rumusan Masalah}\label{rumusan-masalah}

Pertanyaan utama: \textbf{Seberapa stabil dan informatif hasil Regresi
Bayesian pada Yacht Hydrodynamics ketika difokuskan pada reproducibility
\& data quality?} Pertanyaan ini dijabarkan menjadi empat pertanyaan
terukur:

\begin{enumerate}
\def\labelenumi{\arabic{enumi}.}
\tightlist
\item
  Apakah kesimpulan berubah bila seed MCMC atau urutan baris data
  diubah? (ukuran: selisih posterior dalam satuan MCSE)
\item
  Seberapa besar dampak preprocessing (target mentah vs log, bentuk
  fungsi Froude) terhadap kecocokan dan prediksi? (ukuran: ELPD, RMSE,
  cakupan interval prediksi)
\item
  Seberapa optimis validasi acak per baris dibanding validasi per hull?
  (ukuran: selisih ELPD dan RMSE \emph{out-of-sample})
\item
  Seberapa sensitif posterior dan kesimpulan prediktif terhadap pilihan
  prior yang masih masuk akal?
\end{enumerate}

\subsection{1.3 Tujuan Penelitian}\label{tujuan-penelitian}

\begin{enumerate}
\def\labelenumi{\arabic{enumi}.}
\tightlist
\item
  Mengaudit struktur dan kualitas data, termasuk ketergantungan
  observasi dalam satu hull.
\item
  Membangun baseline Regresi Bayesian yang benar dan reproducible, lalu
  memperbaikinya secara bertahap sesuai temuan diagnostik.
\item
  Mengukur stabilitas hasil terhadap seed, urutan data, prior,
  preprocessing, dan skema validasi.
\item
  Mengomunikasikan ketidakpastian estimasi dan prediksi untuk desain
  lambung baru.
\end{enumerate}

\subsection{1.4 Unit Analisis dan
Estimand}\label{unit-analisis-dan-estimand}

{\def\LTcaptype{none} % do not increment counter
\begin{longtable}[]{@{}
  >{\raggedright\arraybackslash}p{(\linewidth - 2\tabcolsep) * \real{0.5000}}
  >{\raggedright\arraybackslash}p{(\linewidth - 2\tabcolsep) * \real{0.5000}}@{}}
\toprule\noalign{}
\begin{minipage}[b]{\linewidth}\raggedright
Elemen
\end{minipage} & \begin{minipage}[b]{\linewidth}\raggedright
Keterangan
\end{minipage} \\
\midrule\noalign{}
\endhead
\bottomrule\noalign{}
\endlastfoot
Unit observasi & Satu uji \emph{towing tank}: satu bentuk lambung pada
satu Froude number \\
Outcome & Hambatan residu per satuan berat displacement
(adimensional) \\
Prediktor & 5 koefisien geometri lambung + Froude number; semua tersedia
saat desain, sehingga tidak ada kebocoran dari outcome \\
Estimand utama & Posterior predictive distribution
\(p(\tilde y \mid \tilde x, y)\) untuk \textbf{lambung baru} \\
Estimand sekunder & Efek 1 SD geometri terhadap resistance; stabilitas
kesimpulan terhadap sumber variasi \\
\end{longtable}
}

\section{BAB 2. TINJAUAN PUSTAKA}\label{bab-2.-tinjauan-pustaka}

\subsection{2.1 Hambatan Residu dan Froude
Number}\label{hambatan-residu-dan-froude-number}

Froude number \(Fr = V/\sqrt{gL}\) adalah kecepatan tak berdimensi
relatif terhadap panjang kapal. Pada \(Fr\) rendah hambatan residu
kecil. Mendekati \(Fr \approx 0.4\), kapal ``terjebak'' di antara
gelombang haluan dan buritannya sendiri sehingga hambatan meningkat
hampir eksponensial (Gerritsma dkk., 1981). Konsekuensinya bagi
pemodelan: hubungan outcome--Froude sangat non-linear, dan variasinya
membesar pada kecepatan tinggi.

\subsection{2.2 Regresi Bayesian}\label{regresi-bayesian}

Inferensi Bayesian memperbarui pengetahuan melalui teorema Bayes:

\[p(\theta\mid y) = \frac{p(y\mid\theta)\,p(\theta)}{p(y)} \propto p(y\mid\theta)\,p(\theta).\]

Untuk regresi Gaussian
\(\mathbf y\mid\beta,\sigma^2 \sim N(\mathbf X\beta, \sigma^2\mathbf I)\)
dengan prior konjugat Normal--Inverse-Gamma
\(\beta\mid\sigma^2\sim N(\mathbf m_0,\sigma^2\mathbf V_0)\) dan
\(\sigma^2\sim IG(a_0,b_0)\), posterior tersedia dalam bentuk tertutup
(Hoff, 2009):

\[\mathbf V_n=(\mathbf V_0^{-1}+\mathbf X^\top\mathbf X)^{-1},\quad \mathbf m_n=\mathbf V_n(\mathbf V_0^{-1}\mathbf m_0+\mathbf X^\top\mathbf y),\quad a_n=a_0+\tfrac n2,\quad b_n=b_0+\tfrac12\left[\mathbf y^\top\mathbf y+\mathbf m_0^\top\mathbf V_0^{-1}\mathbf m_0-\mathbf m_n^\top\mathbf V_n^{-1}\mathbf m_n\right].\]

Prediksi untuk data baru memakai posterior predictive distribution
\(p(\tilde y\mid y)=\int p(\tilde y\mid\theta)\,p(\theta\mid y)\,d\theta\),
yang menggabungkan ketidakpastian parameter dan variasi observasi. Prior
yang baik ditentukan pada skala yang bermakna dan dievaluasi dengan
\textbf{prior predictive check}: \(\theta^{(s)}\sim p(\theta)\), lalu
\(y^{(s)}\sim p(y\mid\theta^{(s)})\) (Gelman dkk., 2020).

\subsection{2.3 Komputasi dan Diagnostik
MCMC}\label{komputasi-dan-diagnostik-mcmc}

Model non-konjugat diestimasi dengan Hamiltonian Monte Carlo (NUTS)
melalui Stan/\texttt{brms} (Bürkner, 2017). Konvergensi dinilai
bersama-sama: \(\hat R\) ter-normalisasi rank ≤ 1.01, \emph{effective
sample size} (bulk \& tail) ≥ 400, tanpa \emph{divergent transitions},
serta trace/rank plot (Vehtari dkk., 2021). Galat Monte Carlo diukur
dengan \(MCSE \approx SD/\sqrt{ESS}\).

\subsection{2.4 Model Checking dan Perbandingan
Model}\label{model-checking-dan-perbandingan-model}

\emph{Posterior predictive check} (PPC) membandingkan statistik data
teramati dengan data replika \(y^{rep}\sim p(y^{rep}\mid y)\) untuk
menemukan aspek data yang gagal direproduksi model. Perbandingan
prediktif memakai \(\text{ELPD}_{LOO}=\sum_i \log p(y_i\mid y_{-i})\)
yang diaproksimasi PSIS-LOO dengan diagnostik Pareto-\(k\) (Vehtari
dkk., 2017). ELPD hanya sah dibandingkan bila semua model mendefinisikan
densitas untuk variabel yang sama. Jika target ditransformasi
\(z=g(y)\), densitas dipindahkan ke skala \(y\) dengan koreksi Jacobian:
\(\log p(y)=\log p(z)+\log|dz/dy|\).

\subsection{2.5 Reproducibility dan Kebocoran
Data}\label{reproducibility-dan-kebocoran-data}

Hasil reproducible berarti dapat direplikasi dari data mentah dan kode
tanpa langkah manual: seed, versi perangkat lunak, dan preprocessing
tercatat. Pada data berkelompok, validasi acak per baris menempatkan
observasi dari kelompok yang sama di data latih dan uji, sehingga model
``mengenal'' kelompok yang diuji (\emph{data leakage}). Estimasi kinerja
menjadi terlalu optimis bila tujuannya memprediksi kelompok baru
(Roberts dkk., 2017).

\section{BAB 3. METODOLOGI}\label{bab-3.-metodologi}

\subsection{3.1 Sumber Data dan
Variabel}\label{sumber-data-dan-variabel}

Data sekunder berasal dari UCI Machine Learning Repository (ID 243,
\emph{Yacht Hydrodynamics}), hasil eksperimen Delft Ship Hydromechanics
Laboratory. File mentah tidak memiliki header, sehingga nama kolom
diberikan mengikuti urutan dokumentasi resmi.

\begingroup\fontsize{8.5}{10.5}\selectfont

\begin{longtable}[t]{>{\raggedleft\arraybackslash}p{1.00cm}>{\raggedright\arraybackslash}p{2.60cm}>{\raggedright\arraybackslash}p{7.70cm}>{\raggedright\arraybackslash}p{5.10cm}}
\caption{\label{tab:tabel_variabel}Tabel 3.1 Variabel Penelitian (semua adimensional)}\\
\toprule
No & Kode & Variabel & Peran\\
\midrule
1 & LCB & Longitudinal position of the center of buoyancy & Prediktor (geometri lambung)\\
2 & Prismatic\_Coef & Prismatic coefficient & Prediktor (geometri lambung)\\
3 & LD\_Ratio & Length-displacement ratio & Prediktor (geometri lambung)\\
4 & BD\_Ratio & Beam-draught ratio & Prediktor (geometri lambung)\\
5 & LB\_Ratio & Length-beam ratio & Prediktor (geometri lambung)\\
\addlinespace
6 & Froude\_No & Froude number & Prediktor (kondisi uji)\\
7 & Resistance & Residuary resistance per unit weight of displacement & Outcome\\
\bottomrule
\end{longtable}
\endgroup{}

\subsection{3.2 Audit Kualitas Data dan
Preprocessing}\label{audit-kualitas-data-dan-preprocessing}

Audit meliputi tipe dan rentang nilai, \emph{missing value}, duplikasi,
struktur hull, distribusi target, dan outlier (IQR). Preprocessing
dilakukan secara eksplisit dalam satu fungsi (\texttt{make\_frames}):
prediktor distandardisasi \(z=(x-\bar x)/s\), target diberi dua versi
(terstandardisasi pada skala mentah, dan \(\log(1+y)\)
terstandardisasi), serta basis B-spline kubik (df = 6) untuk Froude.
\textbf{Statistik standardisasi dan knot spline selalu dihitung dari
data latih saja}, lalu diterapkan ke data uji.

\subsection{3.3 Spesifikasi Model}\label{spesifikasi-model}

Tujuh model dibangun bertahap. Setiap model mengubah satu aspek dari
model sebelumnya, sehingga dampaknya dapat diatribusikan.

\begingroup\fontsize{8.5}{10.5}\selectfont

\begin{longtable}[t]{>{\raggedright\arraybackslash}p{2.00cm}>{\raggedright\arraybackslash}p{6.80cm}>{\raggedright\arraybackslash}p{6.80cm}}
\caption{\label{tab:tabel_model}Tabel 3.2 Model yang Dibandingkan}\\
\toprule
Model & Deskripsi & Perubahan\\
\midrule
M0 & M0 Baseline: Gaussian, raw target (Variant 1) & Baseline: Gaussian, target mentah, linear\\
M1 & M1 Log target, linear (Variant 2) & Target log(1+y) (Varian 2)\\
M2 & M2 Log target + Froude\textasciicircum{}2, \textasciicircum{}3 (Variant 3) & M1 + Froude\textasciicircum{}2, Froude\textasciicircum{}3 (Varian 3)\\
M3 & M3 Log target + Froude spline & Spline B untuk Froude\\
M4 & M4 M3 + heteroskedastic sigma(Froude) & M3 + log sigma = g0 + g1 z\_Froude (heteroskedastik)\\
\addlinespace
M5 & M5 M3 with Student-t likelihood & M3 dengan likelihood Student-t\\
M6 & M6 M4 + hull random intercept (final) & M4 + efek acak hull (dependensi dalam hull)\\
\bottomrule
\end{longtable}
\endgroup{}

Model terlengkap (M6) ditulis sebagai:

\[\ell_{ih}\sim N(\mu_{ih},\sigma_{ih}^2),\qquad \mu_{ih}=\beta_0+\sum_{j\in\text{geom}}\beta_j z_{j,h}+\sum_{k=1}^{6}\alpha_k B_k(z_{F,i})+u_h,\qquad \log\sigma_{ih}=\gamma_0+\gamma_1 z_{F,i},\qquad u_h\sim N(0,\sigma_{hull}^2),\]

dengan \(\ell_{ih}\) adalah \(\log(1+\text{Resistance})\)
terstandardisasi dan \(B_k\) basis B-spline.

\subsection{3.4 Prior}\label{prior}

Semua prior didefinisikan pada skala terstandardisasi dengan
parameterisasi eksplisit, lalu divalidasi dengan prior predictive check
pada \textbf{satuan asli} Resistance.

\begingroup\fontsize{8.5}{10.5}\selectfont

\begin{longtable}[t]{>{\raggedright\arraybackslash}p{3.40cm}>{\raggedright\arraybackslash}p{4.50cm}>{\raggedright\arraybackslash}p{9.00cm}}
\caption{\label{tab:tabel_prior}Tabel 3.3 Prior, Parameterisasi, dan Alasan Substantif}\\
\toprule
Parameter & Prior & Alasan\\
\midrule
Intercept & Normal(0, 1) & Outcome terstandardisasi: di rata-rata prediktor, outcome sekitar 0 plus-minus 2 SD\\
Slope geometri (5) & Normal(0, 0.5) & 22 hull adalah perturbasi kecil di sekitar satu parent hull; regularisasi ridge\\
Slope Froude (F, F\textasciicircum{}2, F\textasciicircum{}3) & Normal(0, 1) & Kecepatan adalah penggerak hambatan yang dominan secara fisik\\
Koefisien spline (6) & Normal(0, 1.5) & Koefisien B-spline = tinggi kurva relatif, rentang log-resistance sekitar 3.4 SD\\
sigma (homoskedastik) & Exponential(rate = 1) & Menjamin sigma > 0; residual SD tidak melebihi SD outcome\\
\addlinespace
log-sigma: gamma0, gamma1 & Normal(-1.5, 0.75); Normal(0, 0.5) & Median sigma = 0.22 SD (95\%: 0.05-0.97); perubahan varians per SD Froude moderat\\
nu (Student-t) & Gamma(shape = 2, rate = 0.1) & Mean 20; mengizinkan ekor berat hingga hampir Normal\\
sd hull & Half-Normal(0, 0.25) & Selisih antar-hull pada geometri sama kurang dari sekitar 0.5 SD\\
\bottomrule
\end{longtable}
\endgroup{}

\subsection{3.5 Komputasi dan Validasi}\label{komputasi-dan-validasi}

Model di-fit dengan \texttt{brms} (Stan, NUTS): 4 chains × 1000 draws
setelah 1000 warm-up, \texttt{adapt\_delta} = 0.95, seed = 2026.
Kebenaran komputasi diverifikasi dengan membandingkan posterior analitik
konjugat, \emph{Gibbs sampler} yang ditulis manual, dan HMC di Stan pada
model dan prior yang sama.

\subsection{3.6 Audit Reproducibility dan Skema
Validasi}\label{audit-reproducibility-dan-skema-validasi}

\begin{enumerate}
\def\labelenumi{\arabic{enumi}.}
\tightlist
\item
  \textbf{Seed:} model final di-fit ulang dengan seed 777; bila ada
  parameter yang ditandai, dilanjutkan dengan audit lima seed (2026,
  777, 999, 31, 123) untuk membandingkan sebaran antar-run dengan MCSE
  yang dilaporkan.
\item
  \textbf{Urutan data:} seed sama, baris data diacak.
\item
  \textbf{Skema validasi:} 5-fold per hull (4--5 hull utuh ditahan tiap
  fold) dibandingkan 5-fold acak per baris. Metrik pada satuan asli:
  RMSE, MAE, cakupan dan lebar interval prediksi 95\%, serta ELPD per
  observasi.
\item
  \textbf{Sensitivitas prior:} semua SD prior M6 dikali 0.5 dan 2.5.
\end{enumerate}

Kriteria stabil: selisih mean posterior \textbar z\textbar{} \textless{}
3 dalam satuan MCSE gabungan \(z=\Delta/\sqrt{MCSE_1^2+MCSE_2^2}\).
Cross-validation memakai konfigurasi ringan (2 chains × 500 draws).

\clearpage

\section{BAB 4. HASIL DAN PEMBAHASAN}\label{bab-4.-hasil-dan-pembahasan}

\subsection{4.1 Audit Kualitas Data}\label{audit-kualitas-data}

\begingroup\fontsize{8.5}{10.5}\selectfont

\begin{longtable}[t]{ll}
\caption{\label{tab:tabel_audit}Tabel 4.1 Ringkasan Audit Kualitas Data}\\
\toprule
Pemeriksaan & Hasil\\
\midrule
Jumlah baris x kolom & 308 x 7\\
Missing value & 0\\
Duplikasi penuh & 0\\
Hull form unik & 22\\
Baris per hull & 14\\
\addlinespace
Level Froude number & 14\\
Skewness target (mentah) & 1.75\\
Skewness target log(1+y) & 0.47\\
Shapiro-Wilk target (p-value) & 4.85e-23\\
Outlier IQR pada target & 42\\
\bottomrule
\end{longtable}
\endgroup{}

Data lengkap tanpa duplikasi, dan desain eksperimennya \textbf{seimbang
sempurna}: 22 hull × 14 Froude number = 308 baris. Temuan terpenting
adalah struktur berkelompok ini: observasi dalam satu hull berbagi
geometri yang sama, sehingga tidak saling independen. Outlier IQR pada
target bukan kesalahan input, melainkan kondisi fisik valid pada
kecepatan tinggi, sehingga \textbf{tidak dihapus}.

\begin{figure}

{\centering \includegraphics[width=0.94\linewidth]{Laporan_Project_Regresi_Bayesian_Yacht_final8_fixed_download_files/figure-latex/fig_froude-1} 

}

\caption{Gambar 4.1 Hubungan Froude Number dan Hambatan Residu per Hull (skala mentah dan log)}\label{fig:fig_froude}
\end{figure}

Setiap hull menunjukkan pola yang sama: datar pada Froude rendah lalu
melonjak tajam setelah \(Fr \approx 0.35\). Pola ini jauh lebih linear
pada skala log. Distribusi target sangat miring ke kanan (skewness
1.75), sedangkan setelah transformasi \(\log(1+y)\) skewness menjadi
0.47. Temuan ini menjadi dasar Varian 2 (target log) dan Varian 3/spline
(bentuk fungsi Froude).

\subsection{4.2 Prior Predictive Check}\label{prior-predictive-check}

\begingroup\fontsize{8.5}{10.5}\selectfont

\begin{longtable}[t]{lrrrrrrr}
\caption{\label{tab:tabel_prior_pred}Tabel 4.2 Prior Predictive pada Satuan Asli Resistance}\\
\toprule
 & share\_negative & share\_below\_minus\_0.5 & share\_gt\_1000 & q01 & q50 & q95 & q99\\
\midrule
M0 (Gaussian, raw) & 0.366 & 0.360 & 0.000 & -75.500 & 10.30 & 6.49e+01 & 9.71e+01\\
M6 (revised prior) & 0.223 & 0.142 & 0.010 & -0.974 & 4.21 & 1.91e+02 & 9.59e+02\\
M6 (initial wider prior) & 0.393 & 0.350 & 0.197 & -1.000 & 3.97 & 5.78e+07 & 6.15e+40\\
Observed data & 0.000 & 0.000 & 0.000 & 0.041 & 3.06 & 4.99e+01 & 5.58e+01\\
\bottomrule
\end{longtable}
\endgroup{}

\begin{figure}

{\centering \includegraphics[width=0.94\linewidth]{Laporan_Project_Regresi_Bayesian_Yacht_final8_fixed_download_files/figure-latex/fig_prior_pred-1} 

}

\caption{Gambar 4.2 Prior Predictive Check: M0 pada skala mentah dan M6 pada skala log}\label{fig:fig_prior_pred}
\end{figure}

Prior predictive check memberi tiga temuan:

\begin{enumerate}
\def\labelenumi{\arabic{enumi}.}
\tightlist
\item
  \textbf{M0} menghasilkan 37\% resistance negatif. Penyebabnya bukan
  prior yang terlalu lebar, melainkan \textbf{likelihood Gaussian pada
  skala mentah yang tidak menghormati support \(y\ge0\)}.
\item
  \textbf{Prior awal M6} tampak \emph{weakly informative} pada skala
  terstandardisasi, tetapi setelah transformasi eksponensial
  menghasilkan resistance absurd (persentil-99 ≈ 6.2e+40). Prior
  kemudian direvisi berdasarkan pemeriksaan ini.
\item
  \textbf{Prior revisi M6} menghasilkan orde besaran yang masuk akal
  (persentil-95 ≈ 191, hanya 1.0\% di atas 1000). Nilai sedikit negatif
  (−1 \textless{} R \textless{} 0) masih muncul karena \(\log(1+y)\)
  mengizinkan \(y>-1\); ini dicatat sebagai keterbatasan.
\end{enumerate}

\textbf{Koreksi versi awal (temuan audit).} Versi awal analisis memasang
prior \(N(0,1)\) langsung pada koefisien prediktor skala mentah. Prior
itu menyatakan bahwa kenaikan Froude dari minimum ke maksimum hanya
mengubah resistance sekitar ±0.64, padahal data naik dari ≈0 ke ≈62.
Hasilnya, koefisien Froude tertarik menjadi 1.69 dibanding OLS 121.67,
\textbf{tanpa satu pun peringatan konvergensi}. Kasus ini menunjukkan
mengapa standardisasi eksplisit dan prior predictive check wajib
dilakukan.

\subsection{4.3 Validasi Komputasi}\label{validasi-komputasi}

\begingroup\fontsize{8.5}{10.5}\selectfont

\begin{longtable}[t]{lrrrrrrr}
\caption{\label{tab:tabel_komputasi}Tabel 4.3 OLS vs Posterior Analitik vs Gibbs Manual vs HMC (Stan), Prior NIG Identik}\\
\toprule
 & OLS & Analytic\_mean & Analytic\_SD & Gibbs\_mean & HMC\_mean & HMC\_MCSE & HMC\_minus\_Analytic\_in\_MCSE\\
\midrule
Intercept & 0.0000 & 0.0000 & 0.0336 & 0.0001 & 0.0000 & 0.0004 & 0.0534\\
z\_LCB & 0.0193 & 0.0192 & 0.0336 & 0.0193 & 0.0189 & 0.0004 & -0.6783\\
z\_Prismatic\_Coef & -0.0099 & -0.0151 & 0.0596 & -0.0147 & -0.0162 & 0.0011 & -1.0175\\
z\_LD\_Ratio & 0.0707 & 0.0492 & 0.1975 & 0.0497 & 0.0459 & 0.0037 & -0.8781\\
z\_BD\_Ratio & -0.0638 & -0.0458 & 0.1670 & -0.0458 & -0.0430 & 0.0032 & 0.8771\\
\addlinespace
z\_LB\_Ratio & -0.0739 & -0.0528 & 0.1941 & -0.0533 & -0.0497 & 0.0037 & 0.8465\\
z\_Froude\_No & 0.8101 & 0.8075 & 0.0336 & 0.8075 & 0.8077 & 0.0003 & 0.5910\\
sigma2 & NA & 0.3488 & NA & 0.3488 & 0.3485 & NA & NA\\
\bottomrule
\end{longtable}
\endgroup{}

Keempat pendekatan memberikan hasil praktis identik. Selisih terbesar
HMC terhadap solusi analitik adalah 1.02 MCSE, konsisten dengan galat
Monte Carlo. Log-likelihood yang ditulis manual (-1,115.9416) identik
dengan log-likelihood Stan (-1,115.9416). Jadi komputasi benar;
ketidakstabilan yang mungkin muncul bukan berasal dari sampler.

\subsection{4.4 Baseline M0 dan Diagnosis
Kegagalannya}\label{baseline-m0-dan-diagnosis-kegagalannya}

\begingroup\fontsize{8.5}{10.5}\selectfont

\begin{longtable}[t]{lrrrrrrrr}
\caption{\label{tab:tabel_m0}Tabel 4.4 Posterior Koefisien M0 dalam Satuan Asli (resistance per 1 unit prediktor)}\\
\toprule
 & mean & sd & eti\_2.5 & eti\_97.5 & hdi\_low & hdi\_high & P\_gt\_0 & OLS\_raw\_units\\
\midrule
LCB & 0.19 & 0.34 & -0.47 & 0.87 & -0.50 & 0.83 & 0.71 & 0.19\\
Prismatic\_Coef & -9.65 & 37.35 & -80.20 & 64.06 & -82.12 & 61.18 & 0.40 & -6.42\\
LD\_Ratio & 2.78 & 10.95 & -18.23 & 24.58 & -18.24 & 24.58 & 0.60 & 4.23\\
BD\_Ratio & -1.21 & 4.27 & -9.81 & 7.12 & -10.11 & 6.79 & 0.39 & -1.77\\
LB\_Ratio & -3.06 & 11.03 & -25.37 & 18.47 & -26.15 & 17.59 & 0.39 & -4.52\\
\addlinespace
Froude\_No & 121.52 & 5.24 & 111.34 & 132.03 & 110.87 & 131.44 & 1.00 & 121.67\\
\bottomrule
\end{longtable}
\endgroup{}

Slope Froude pada M0 = 121.5 dengan 95\% \emph{credible interval}
{[}111.3, 132.0{]}, dekat dengan OLS. Interpretasi Bayesian-nya:
\(P(\beta_{Fr}\in[\text{batas}]\mid y)=0.95\), bersyarat pada model dan
prior --- berbeda dari \emph{confidence interval} frequentist. Namun
angka ini \textbf{belum layak diinterpretasikan} karena modelnya gagal
pada model checking:

\begin{figure}

{\centering \includegraphics[width=0.94\linewidth]{Laporan_Project_Regresi_Bayesian_Yacht_final8_fixed_download_files/figure-latex/fig_ppc_m0-1} 

}

\caption{Gambar 4.3 Posterior Predictive Check Model M0 (garis gelap = statistik data teramati)}\label{fig:fig_ppc_m0}
\end{figure}

\begingroup\fontsize{8.5}{10.5}\selectfont

\begin{longtable}[t]{lrrrrr}
\caption{\label{tab:tabel_ppc_m0}Tabel 4.5 Discrepancy Statistics PPC Model M0}\\
\toprule
statistic & observed & rep\_median & rep\_2.5 & rep\_97.5 & P\_rep\_ge\_obs\\
\midrule
mean & 10.495 & 10.511 & 9.186 & 11.931 & 0.510\\
sd & 15.160 & 15.216 & 13.828 & 16.496 & 0.534\\
max & 62.420 & 49.087 & 42.218 & 58.040 & 0.004\\
skewness & 1.748 & -0.003 & -0.197 & 0.196 & 0.000\\
share < 0 & 0.000 & 0.266 & 0.221 & 0.312 & 1.000\\
\bottomrule
\end{longtable}
\endgroup{}

\begin{figure}

{\centering \includegraphics[width=0.94\linewidth]{Laporan_Project_Regresi_Bayesian_Yacht_final8_fixed_download_files/figure-latex/fig_resid_m0-1} 

}

\caption{Gambar 4.4 Diagnostik Residual Model M0}\label{fig:fig_resid_m0}
\end{figure}

M0 mereplikasi \emph{mean} dan \emph{SD}, tetapi gagal pada aspek yang
paling penting secara fisik: skewness dan maksimum tidak tereplikasi,
sekitar 27\% replikasi bernilai negatif, residual membentuk pola U
terhadap Froude (\textbf{struktur mean salah}), dan sebarannya melebar
pada Froude tinggi (\textbf{heteroskedastisitas}). Sesuai prinsip
``ganti likelihood, bukan memaksa data'', model diperbaiki secara
bertahap (M1--M6).

\subsection{4.5 Model Remedi dan Diagnostik
Konvergensi}\label{model-remedi-dan-diagnostik-konvergensi}

\begingroup\fontsize{8.5}{10.5}\selectfont

\begin{longtable}[t]{lrrrrrrl}
\caption{\label{tab:tabel_diag}Tabel 4.6 Diagnostik MCMC Seluruh Model (waktu sampling dalam detik)}\\
\toprule
model & sampling\_s & max\_rhat & min\_ess\_bulk & min\_ess\_tail & max\_mcse\_over\_sd & divergences & verdict\\
\midrule
M0 & 0.8 & 1.004 & 1175.753 & 1697.216 & 0.029 & 0 & OK\\
M1 & 0.9 & 1.006 & 1204.603 & 1422.063 & 0.029 & 0 & OK\\
M2 & 1.0 & 1.002 & 1092.401 & 1543.522 & 0.030 & 0 & OK\\
M3 & 1.4 & 1.003 & 1315.544 & 1640.096 & 0.028 & 0 & OK\\
M4 & 18.7 & 1.007 & 1400.677 & 2042.288 & 0.027 & 0 & OK\\
\addlinespace
M5 & 10.2 & 1.003 & 1389.435 & 1927.189 & 0.027 & 0 & OK\\
M6 & 149.7 & 1.004 & 1036.491 & 1643.566 & 0.031 & 0 & OK\\
\bottomrule
\end{longtable}
\endgroup{}

Seluruh model memenuhi kriteria konvergensi (R̂ ≤ 1.01, ESS ≥ 400, tanpa
divergensi).

\begin{figure}

{\centering \includegraphics[width=0.94\linewidth]{Laporan_Project_Regresi_Bayesian_Yacht_final8_fixed_download_files/figure-latex/fig_trace-1} 

}

\caption{Gambar 4.5 Trace Plot dan Rank Plot Model M6}\label{fig:fig_trace}
\end{figure}

\begingroup\fontsize{8.5}{10.5}\selectfont

\begin{longtable}[t]{lllr}
\caption{\label{tab:tabel_vif}Tabel 4.7 Variance Inflation Factor}\\
\toprule
 & Kelompok & Prediktor & VIF\\
\midrule
LCB & Geometri + Froude (M0, M1) & LCB & 1.00\\
Prismatic\_Coef & Geometri + Froude (M0, M1) & Prismatic\_Coef & 4.05\\
LD\_Ratio & Geometri + Froude (M0, M1) & LD\_Ratio & 49.14\\
BD\_Ratio & Geometri + Froude (M0, M1) & BD\_Ratio & 35.03\\
LB\_Ratio & Geometri + Froude (M0, M1) & LB\_Ratio & 47.43\\
\addlinespace
Froude\_No & Geometri + Froude (M0, M1) & Froude\_No & 1.00\\
Froude & Suku polinomial Froude (M2) & Froude & 6.48\\
Froude\textasciicircum{}2 & Suku polinomial Froude (M2) & Froude\textasciicircum{}2 & 1.00\\
Froude\textasciicircum{}3 & Suku polinomial Froude (M2) & Froude\textasciicircum{}3 & 6.48\\
\bottomrule
\end{longtable}
\endgroup{}

Multikolinearitas terbesar terdapat pada rasio geometri LD, BD, dan LB
(VIF 35--49). Hal ini wajar karena ketiganya dibentuk dari dimensi yang
sama (panjang, lebar, sarat, volume). Prior \(N(0, 0.5)\) bekerja
sebagai regularisasi ridge yang menstabilkan estimasi, tetapi
\textbf{tidak dapat menciptakan informasi} untuk memisahkan efek
masing-masing rasio (Gambar 4.6).

\begin{figure}

{\centering \includegraphics[width=0.94\linewidth]{Laporan_Project_Regresi_Bayesian_Yacht_final8_fixed_download_files/figure-latex/fig_cor-1} 

}

\caption{Gambar 4.6 Korelasi Posterior Antar-Koefisien (M2, Varian 3)}\label{fig:fig_cor}
\end{figure}

\subsection{4.6 Perbandingan Model: PSIS-LOO dan
WAIC}\label{perbandingan-model-psis-loo-dan-waic}

Karena M0 memodelkan skala mentah sedangkan M1--M6 memodelkan
\(\log(1+y)\), log-likelihood setiap model dipindahkan ke satuan asli
Resistance dengan koreksi Jacobian, yaitu \(-\log s_y\) untuk M0 dan
\(-\log s_\ell-\log(1+y)\) untuk M1--M6. Tanpa koreksi ini perbandingan
\textbf{tidak sah}.

\begingroup\fontsize{8.5}{10.5}\selectfont

\begin{longtable}[t]{lrrrrrlrr}
\caption{\label{tab:tabel_loo}Tabel 4.8 Perbandingan Model dengan PSIS-LOO dan WAIC (satuan asli, koreksi Jacobian)}\\
\toprule
model & elpd\_loo & se\_elpd\_loo & p\_loo & elpd\_diff & se\_diff & elpd\_waic & n\_pareto\_k\_gt\_0.7 & diff\_over\_se\\
\midrule
NULL & -122.31 & 32.51 & 28.53 & 0.00 & 0.00 & NULL & NA & NA\\
NULL & -154.48 & 27.88 & 13.31 & -32.17 & 9.95 & NULL & NA & -3.23\\
NULL & -161.98 & 28.17 & 13.35 & -39.67 & 11.15 & NULL & NA & -3.56\\
NULL & -163.55 & 28.00 & 12.92 & -41.25 & 11.27 & NULL & NA & -3.66\\
NULL & -187.74 & 28.97 & 10.24 & -65.44 & 13.57 & NULL & NA & -4.82\\
\addlinespace
NULL & -509.54 & 24.92 & 7.64 & -387.23 & 17.10 & NULL & NA & -22.64\\
NULL & -1117.33 & 14.83 & 8.77 & -995.02 & 28.78 & NULL & NA & -34.58\\
\bottomrule
\end{longtable}
\endgroup{}

\begin{figure}

{\centering \includegraphics[width=0.94\linewidth]{Laporan_Project_Regresi_Bayesian_Yacht_final8_fixed_download_files/figure-latex/fig_loo-1} 

}

\caption{Gambar 4.7 ELPD-LOO Seluruh Model (plus-minus 1 SE)}\label{fig:fig_loo}
\end{figure}

Perbaikan terbesar datang dari \textbf{preprocessing target dan bentuk
fungsi Froude} (M0 → M1 → M2/M3), disusul pemodelan varians (M4) dan
efek hull (M6). Namun perlu kehati-hatian: PSIS-LOO meninggalkan
\emph{satu baris}, sementara 13 baris lain dari hull yang sama tetap
membantu mengestimasi efek hull tersebut. Keunggulan M6 di LOO karena
itu harus diuji ulang dengan validasi per hull (Subbab 4.10).

\subsection{4.7 Model Checking dan Posterior Model Final
(M6)}\label{model-checking-dan-posterior-model-final-m6}

\begin{figure}

{\centering \includegraphics[width=0.94\linewidth]{Laporan_Project_Regresi_Bayesian_Yacht_final8_fixed_download_files/figure-latex/fig_ppc_m6-1} 

}

\caption{Gambar 4.8 Posterior Predictive Check Model M6}\label{fig:fig_ppc_m6}
\end{figure}

\begingroup\fontsize{8.5}{10.5}\selectfont

\begin{longtable}[t]{lrrrrrrr}
\caption{\label{tab:tabel_ppc_cmp}Tabel 4.9 Perbandingan Discrepancy Statistics M0 dan M6}\\
\toprule
Statistik & Observed & M0\_rep\_median & M0\_P\_rep\_ge\_obs & M6\_rep\_median & M6\_rep\_2.5 & M6\_rep\_97.5 & M6\_P\_rep\_ge\_obs\\
\midrule
mean & 10.495 & 10.511 & 0.510 & 10.550 & 10.203 & 10.922 & 0.613\\
sd & 15.160 & 15.216 & 0.534 & 15.416 & 14.552 & 16.338 & 0.712\\
max & 62.420 & 49.087 & 0.004 & 68.371 & 58.736 & 82.847 & 0.881\\
skewness & 1.748 & -0.003 & 0.000 & 1.814 & 1.683 & 1.980 & 0.814\\
share < 0 & 0.000 & 0.266 & 1.000 & 0.003 & 0.000 & 0.010 & 1.000\\
\bottomrule
\end{longtable}
\endgroup{}

Pada M6 seluruh statistik teramati berada di dalam distribusi replikasi.
Sisa kecil adalah 0.3\% replikasi bernilai sedikit negatif, akibat sifat
transformasi \(\log(1+y)\).

\begin{figure}

{\centering \includegraphics[width=0.94\linewidth]{Laporan_Project_Regresi_Bayesian_Yacht_final8_fixed_download_files/figure-latex/fig_resid_m6-1} 

}

\caption{Gambar 4.9 Residual Terstandardisasi M6 (residual dibagi sigma yang dimodelkan)}\label{fig:fig_resid_m6}
\end{figure}

Transformasi ke besaran yang bermakna dihitung langsung dari setiap draw
posterior: \(+1\) SD geometri mengubah \((1+\text{Resistance})\) sebesar
\(100[\exp(\beta_j s_\ell)-1]\%\).

\begingroup\fontsize{8.5}{10.5}\selectfont

\begin{longtable}[t]{lrrrrrrr}
\caption{\label{tab:tabel_eff}Tabel 4.10 Posterior Efek pada Model M6 (median, 95\% ETI dan HDI, posterior probability)}\\
\toprule
 & median & eti\_2.5 & eti\_97.5 & hdi\_low & hdi\_high & P\_gt\_0 & P\_abs\_gt\_1pct\\
\midrule
LCB: \% change per +1 SD & 1.470 & -0.719 & 3.536 & -0.719 & 3.535 & 0.924 & 0.697\\
Prismatic\_Coef: \% change per +1 SD & 0.103 & -4.110 & 4.377 & -4.179 & 4.240 & 0.524 & 0.617\\
LD\_Ratio: \% change per +1 SD & 4.445 & -10.530 & 20.728 & -11.148 & 19.860 & 0.740 & 0.903\\
BD\_Ratio: \% change per +1 SD & -0.365 & -11.969 & 13.123 & -12.225 & 12.740 & 0.477 & 0.855\\
LB\_Ratio: \% change per +1 SD & -3.779 & -16.688 & 11.946 & -17.346 & 10.805 & 0.289 & 0.901\\
\addlinespace
Residual SD multiplier per +1 SD Froude & 1.598 & 1.388 & 1.848 & 1.370 & 1.828 & 1.000 & 1.000\\
Between-hull variation (\%) & 4.611 & 3.039 & 7.494 & 2.740 & 6.918 & 1.000 & 1.000\\
\bottomrule
\end{longtable}
\endgroup{}

\begin{figure}

{\centering \includegraphics[width=0.94\linewidth]{Laporan_Project_Regresi_Bayesian_Yacht_final8_fixed_download_files/figure-latex/fig_eff-1} 

}

\caption{Gambar 4.10 Efek Geometri Lambung pada Model M6}\label{fig:fig_eff}
\end{figure}

Setelah bentuk fungsi Froude, varians, dan variasi antar-hull
dimodelkan, efek per SD geometri hanya beberapa persen dan
\textbf{arahnya tidak pasti}: \(P(\text{efek}>0\mid y)\) berkisar
0.29--0.92. Sebaliknya, dua besaran diestimasi dengan sangat pasti. SD
residual naik 1.60× per SD Froude (95\% CrI 1.39--1.85; Gambar 4.11).
Variasi antar-hull yang tidak dijelaskan geometri adalah 4.6\%
(3.0--7.5\%).

\begin{figure}

{\centering \includegraphics[width=0.94\linewidth]{Laporan_Project_Regresi_Bayesian_Yacht_final8_fixed_download_files/figure-latex/fig_sigma-1} 

}

\caption{Gambar 4.11 Heteroskedastisitas yang Dimodelkan: SD Residual Meningkat dengan Froude Number}\label{fig:fig_sigma}
\end{figure}

\subsection{4.8 Prediksi untuk Desain Lambung
Baru}\label{prediksi-untuk-desain-lambung-baru}

Skenario: lambung \textbf{baru} dengan geometri di median ke-22 desain,
diuji pada 14 Froude number eksperimen (tanpa ekstrapolasi). Efek hull
baru diambil dari distribusi populasi
\(u_{new}\sim N(0,\sigma_{hull}^2)\).

\begin{figure}

{\centering \includegraphics[width=0.94\linewidth]{Laporan_Project_Regresi_Bayesian_Yacht_final8_fixed_download_files/figure-latex/fig_pred-1} 

}

\caption{Gambar 4.12 Posterior Predictive untuk Lambung Baru: Credible Band vs Predictive Band}\label{fig:fig_pred}
\end{figure}

\begingroup\fontsize{8.5}{10.5}\selectfont

\begin{longtable}[t]{rrrrrr}
\caption{\label{tab:tabel_pred}Tabel 4.11 Prediksi Lambung Baru: Median, Interval Prediksi 95\%, dan Probabilitas Eksedans}\\
\toprule
Froude & pred\_median & PI\_2.5 & PI\_97.5 & P\_R\_gt\_5 & P\_R\_gt\_20\\
\midrule
0.125 & 0.113 & -0.008 & 0.255 & 0.000 & 0.000\\
0.150 & 0.265 & 0.121 & 0.430 & 0.000 & 0.000\\
0.175 & 0.488 & 0.313 & 0.698 & 0.000 & 0.000\\
0.200 & 0.809 & 0.585 & 1.066 & 0.000 & 0.000\\
0.225 & 1.243 & 0.949 & 1.586 & 0.000 & 0.000\\
\addlinespace
0.250 & 1.851 & 1.455 & 2.290 & 0.000 & 0.000\\
0.275 & 2.682 & 2.152 & 3.315 & 0.000 & 0.000\\
0.300 & 3.829 & 3.068 & 4.717 & 0.005 & 0.000\\
0.325 & 5.484 & 4.436 & 6.773 & 0.803 & 0.000\\
0.350 & 8.169 & 6.478 & 10.231 & 1.000 & 0.000\\
\addlinespace
0.375 & 12.992 & 10.233 & 16.320 & 1.000 & 0.000\\
0.400 & 21.866 & 16.913 & 28.100 & 1.000 & 0.764\\
0.425 & 35.852 & 27.247 & 47.752 & 1.000 & 1.000\\
0.450 & 51.852 & 37.699 & 69.270 & 1.000 & 1.000\\
\bottomrule
\end{longtable}
\endgroup{}

Probabilitas eksedans seperti \(P(R>20\mid y)\) = 0.76 pada
\(Fr = 0.40\) dapat langsung dipakai untuk keputusan desain (misalnya
kebutuhan daya propulsi). Band prediksi jauh lebih lebar daripada
\emph{credible band}: perbedaan antara \textbf{ketidakpastian estimasi}
dan \textbf{ketidakpastian prediksi} adalah pesan utama dalam
mengomunikasikan uncertainty.

\subsection{4.9 Analisis Sensitivitas}\label{analisis-sensitivitas}

\subsubsection{Kompetisi Prior dan Data}\label{kompetisi-prior-dan-data}

Data dan likelihood dibuat tetap, hanya prior slope Froude yang diubah.
Posterior konjugat dihitung untuk tiga ukuran data berbasis hull.

\begin{figure}

{\centering \includegraphics[width=0.94\linewidth]{Laporan_Project_Regresi_Bayesian_Yacht_final8_fixed_download_files/figure-latex/fig_lab-1} 

}

\caption{Gambar 4.13 Kompetisi Prior-Data pada Berbagai Ukuran Data}\label{fig:fig_lab}
\end{figure}

\begingroup\fontsize{8.5}{10.5}\selectfont

\begin{longtable}[t]{lllrrrr}
\caption{\label{tab:tabel_lab}Tabel 4.12 Kompetisi Prior-Data (prior pull = pangsa presisi prior dalam presisi posterior)}\\
\toprule
 & data & prior & post\_mean & post\_sd & data\_only\_slope & prior\_pull\\
\midrule
z\_Froude\_No & 1 hull (n=14) & Weak N(0, 1) & 0.739 & 0.156 & 0.758 & 0.024\\
z\_Froude\_No1 & 1 hull (n=14) & Strong-correct N(0.81, 0.02) & 0.809 & 0.020 & 0.758 & 0.984\\
z\_Froude\_No2 & 1 hull (n=14) & Strong-wrong N(0.40, 0.02) & 0.406 & 0.020 & 0.758 & 0.984\\
z\_Froude\_No3 & 3 hulls (n=42) & Weak N(0, 1) & 0.772 & 0.091 & 0.778 & 0.008\\
z\_Froude\_No11 & 3 hulls (n=42) & Strong-correct N(0.81, 0.02) & 0.809 & 0.020 & 0.778 & 0.954\\
\addlinespace
z\_Froude\_No21 & 3 hulls (n=42) & Strong-wrong N(0.40, 0.02) & 0.417 & 0.020 & 0.778 & 0.954\\
z\_Froude\_No4 & 22 hulls (n=308) & Weak N(0, 1) & 0.809 & 0.034 & 0.810 & 0.001\\
z\_Froude\_No12 & 22 hulls (n=308) & Strong-correct N(0.81, 0.02) & 0.810 & 0.017 & 0.810 & 0.739\\
z\_Froude\_No22 & 22 hulls (n=308) & Strong-wrong N(0.40, 0.02) & 0.507 & 0.017 & 0.810 & 0.739\\
\bottomrule
\end{longtable}
\endgroup{}

Prior lemah praktis menghilang begitu ada ≥ 3 hull. Prior kuat yang
\textbf{salah} dengan SD 0.02 tetap menarik posterior bahkan pada n =
308, karena SD prior itu lebih kecil daripada ketidakpastian yang
diberikan data --- prior ekstrem yang sangat terkonsentrasi dapat
mendominasi lebih lama.

\subsubsection{Sensitivitas Prior pada Model
Final}\label{sensitivitas-prior-pada-model-final}

\begin{figure}

{\centering \includegraphics[width=0.94\linewidth]{Laporan_Project_Regresi_Bayesian_Yacht_final8_fixed_download_files/figure-latex/fig_sens-1} 

}

\caption{Gambar 4.14 Sensitivitas Posterior M6 terhadap Skala Prior}\label{fig:fig_sens}
\end{figure}

\begingroup\fontsize{8.5}{10.5}\selectfont

\begin{longtable}[t]{lrr}
\caption{\label{tab:tabel_sens}Tabel 4.13 Kesimpulan Prediktif di Bawah Tiga Skala Prior}\\
\toprule
prior & P\_R20\_Fr040 & ELPD\_LOO\\
\midrule
Tighter (x0.5) & 0.761 & -123.082\\
Base (x1) & 0.764 & -122.307\\
Wider (x2.5) & 0.763 & -122.475\\
\bottomrule
\end{longtable}
\endgroup{}

Mengubah semua SD prior dari ×0.5 hingga ×2.5 praktis tidak mengubah
posterior maupun kesimpulan prediktif (\(P(R>20\mid Fr=0.40)\) berkisar
0.76--0.76). Untuk sensitivitas likelihood, posterior derajat bebas
Student-\(t\) (M5) memiliki median 14.0, dan ELPD M5 − M3 = 1.57.
Artinya, setelah transformasi log, ekor berat tidak diperlukan;
perbaikan berasal dari struktur mean dan varians.

\subsection{4.10 Audit Reproducibility}\label{audit-reproducibility}

\subsubsection{Seed dan Urutan Data}\label{seed-dan-urutan-data}

\begin{figure}

{\centering \includegraphics[width=0.94\linewidth]{Laporan_Project_Regresi_Bayesian_Yacht_final8_fixed_download_files/figure-latex/fig_repro-1} 

}

\caption{Gambar 4.15 Selisih Mean Posterior terhadap Run Referensi (satuan MCSE gabungan)}\label{fig:fig_repro}
\end{figure}

Selisih maksimum adalah 1.26 MCSE (seed) dan 1.84 MCSE (urutan data).
Seluruh parameter memenuhi kriteria \textbar z\textbar{} \textless{} 3,
sehingga perbedaan hanya berasal dari Monte Carlo error.

\begingroup\fontsize{8.5}{10.5}\selectfont

\begin{longtable}[t]{lrrrrr}
\caption{\label{tab:tabel_multiseed}Tabel 4.14 Audit Lima Seed: Sebaran Mean Posterior Antar-Run vs MCSE yang Dilaporkan}\\
\toprule
param & min\_mean\_5\_seeds & max\_mean\_5\_seeds & between\_run\_sd & mean\_reported\_mcse & ratio\_sd\_to\_mcse\\
\midrule
b\_F1 & 0.08693 & 0.08769 & 0.00032 & 0.00022 & 1.42009\\
b\_Intercept & -1.28009 & -1.27956 & 0.00023 & 0.00020 & 1.14608\\
b\_sigma\_Intercept & -2.94556 & -2.94405 & 0.00073 & 0.00072 & 1.01123\\
b\_z\_LCB & 0.01185 & 0.01230 & 0.00018 & 0.00019 & 0.95647\\
b\_z\_Prismatic\_Coef & -0.00012 & 0.00077 & 0.00039 & 0.00042 & 0.93260\\
\addlinespace
b\_z\_BD\_Ratio & -0.00307 & -0.00030 & 0.00107 & 0.00126 & 0.84877\\
b\_F5 & 2.93221 & 2.93335 & 0.00045 & 0.00053 & 0.83449\\
b\_z\_LB\_Ratio & -0.03061 & -0.02746 & 0.00122 & 0.00147 & 0.82661\\
b\_z\_LD\_Ratio & 0.03232 & 0.03525 & 0.00120 & 0.00149 & 0.80373\\
b\_F6 & 3.14387 & 3.14448 & 0.00023 & 0.00036 & 0.63167\\
\addlinespace
b\_sigma\_z\_Froude\_No & 0.46702 & 0.46966 & 0.00095 & 0.00157 & 0.60436\\
b\_F4 & 1.85962 & 1.86030 & 0.00027 & 0.00046 & 0.57985\\
sd\_hull\_\_Intercept & 0.03760 & 0.03793 & 0.00014 & 0.00027 & 0.50606\\
b\_F3 & 1.05966 & 1.06002 & 0.00014 & 0.00035 & 0.39325\\
b\_F2 & 0.38851 & 0.38868 & 0.00007 & 0.00027 & 0.26574\\
\bottomrule
\end{longtable}
\endgroup{}

\begingroup\fontsize{8.5}{10.5}\selectfont

\begin{longtable}[t]{lrr}
\caption{\label{tab:tabel_seedconcl}Tabel 4.15 Kesimpulan Substantif di Bawah Lima Seed}\\
\toprule
seed & between\_hull\_pct & P\_R20\_Fr040\\
\midrule
2026 & 4.611 & 0.764\\
777 & 4.609 & 0.757\\
999 & 4.560 & 0.757\\
31 & 4.567 & 0.765\\
123 & 4.612 & 0.757\\
\bottomrule
\end{longtable}
\endgroup{}

Rasio sebaran antar-run terhadap MCSE bernilai sekitar 1 untuk sebagian
besar parameter (median 0.83). Untuk SD efek hull, rasionya 0.51:
\textbf{MCSE meremehkan variabilitas nyata} parameter skala hierarkis,
dan run referensi (seed 2026) kebetulan berada di sisi bawah. Secara
substantif perbedaannya kecil: variasi antar-hull berkisar 4.56--4.61\%
dan \(P(R>20\mid Fr=0.40)\) berkisar 0.757--0.765 antar seed.
Pelajarannya: untuk parameter skala hierarkis, satu run dan MCSE-nya
belum cukup sebagai bukti reproducibility. Gunakan beberapa seed (atau
lebih banyak draws), dan nilai stabilitas pada besaran yang menjadi
dasar keputusan.

\subsubsection{Validasi Per Hull vs Acak Per Baris dan Dampak
Preprocessing}\label{validasi-per-hull-vs-acak-per-baris-dan-dampak-preprocessing}

\begingroup\fontsize{8.5}{10.5}\selectfont

\begin{longtable}[t]{llrrrrrrr}
\caption{\label{tab:tabel_cv}Tabel 4.16 Hasil Cross-Validation 5-Fold (satuan asli Resistance)}\\
\toprule
model & scheme & RMSE & MAE & coverage\_95 & PI\_width & ELPD\_per\_obs & ELPD\_total & max\_rhat\_folds\\
\midrule
M0 & grouped & 8.866 & 7.131 & 0.942 & 35.993 & -3.603 & -1109.732 & 1.022\\
M1 & grouped & 4.942 & 2.512 & 0.971 & 9.973 & -1.640 & -504.991 & 1.024\\
M2 & grouped & 2.137 & 0.953 & 0.942 & 3.852 & -0.629 & -193.783 & 1.020\\
M3 & grouped & 1.866 & 0.841 & 0.938 & 3.537 & -0.567 & -174.570 & 1.032\\
M4 & grouped & 2.014 & 0.901 & 0.919 & 4.286 & -0.578 & -177.927 & 1.013\\
\addlinespace
M6 & grouped & 2.091 & 0.934 & 0.987 & 6.096 & -0.628 & -193.297 & 1.009\\
M4 & row-random & 1.870 & 0.848 & 0.938 & 4.188 & -0.523 & -160.971 & 1.011\\
M6 & row-random & 2.063 & 0.927 & 0.938 & 5.075 & -0.422 & -129.881 & 1.018\\
\bottomrule
\end{longtable}
\endgroup{}

\begin{figure}

{\centering \includegraphics[width=0.94\linewidth]{Laporan_Project_Regresi_Bayesian_Yacht_final8_fixed_download_files/figure-latex/fig_cv-1} 

}

\caption{Gambar 4.16 Dampak Preprocessing dan Model pada Prediksi Hull Baru (leave-hull-out)}\label{fig:fig_cv}
\end{figure}

\begingroup\fontsize{8.5}{10.5}\selectfont

\begin{longtable}[t]{lrrrrr}
\caption{\label{tab:tabel_paired}Tabel 4.17 Perbandingan Berpasangan Leave-Hull-Out terhadap Model Terbaik (M3); SE per cluster hull}\\
\toprule
model & ELPD\_total & elpd\_diff & se\_diff\_obs & se\_diff\_hull & diff\_over\_se\_hull\\
\midrule
M3 & -174.57 & 0.00 & 0.00 & 0.00 & NA\\
M4 & -177.93 & -3.36 & 5.19 & 9.12 & -0.37\\
M6 & -193.30 & -18.73 & 8.45 & 18.14 & -1.03\\
M2 & -193.78 & -19.21 & 7.93 & 7.79 & -2.47\\
M1 & -504.99 & -330.42 & 13.84 & 26.82 & -12.32\\
\addlinespace
M0 & -1109.73 & -935.16 & 24.64 & 27.31 & -34.24\\
\bottomrule
\end{longtable}
\endgroup{}

\begingroup\fontsize{8.5}{10.5}\selectfont

\begin{longtable}[t]{lrrrrr}
\caption{\label{tab:tabel_optimism}Tabel 4.18 Optimisme Validasi Acak Per Baris dibanding Per Hull}\\
\toprule
model & ELPD\_optimism & se & RMSE\_grouped & RMSE\_row & optimism\_over\_se\\
\midrule
M4 & 16.96 & 4.74 & 2.01 & 1.87 & 3.58\\
M6 & 63.42 & 11.69 & 2.09 & 2.06 & 5.42\\
\bottomrule
\end{longtable}
\endgroup{}

\begin{figure}

{\centering \includegraphics[width=0.94\linewidth]{Laporan_Project_Regresi_Bayesian_Yacht_final8_fixed_download_files/figure-latex/fig_leak-1} 

}

\caption{Gambar 4.17 Kebocoran Data: Validasi Acak Per Baris Terlalu Optimis untuk Hull Baru}\label{fig:fig_leak}
\end{figure}

Temuan kunci audit validasi:

\begin{enumerate}
\def\labelenumi{\arabic{enumi}.}
\tightlist
\item
  \textbf{Preprocessing tidak stabil dan sangat menentukan.} RMSE per
  hull turun dari 8.87 (M0) ke 1.87 (M3); ELPD per observasi naik dari
  -3.60 ke -0.57. Cakupan M0 (0.94) tampak wajar hanya karena
  intervalnya sangat lebar (36.0) dan mencakup resistance negatif.
\item
  \textbf{Skema validasi mengubah model yang terpilih.} PSIS-LOO per
  baris memilih 1, sedangkan validasi per hull memilih M3. Selisih antar
  model berbasis spline (M3, M4, M6) untuk hull baru ≤ 1.0 SE (per
  cluster hull), sehingga ketiganya praktis setara.
\item
  \textbf{Kebocoran data terukur.} Validasi acak per baris
  melebih-lebihkan ELPD M6 sebesar 63.4 (SE 11.7), jauh lebih besar
  daripada M4 (17.0, SE 4.7), karena M6 ``mengingat'' efek tiap hull.
  Kebocoran terutama memengaruhi \textbf{lebar interval}, bukan titik
  prediksi.
\end{enumerate}

\begingroup\fontsize{8.5}{10.5}\selectfont

\begin{longtable}[t]{ll}
\caption{\label{tab:tabel_check}Tabel 4.19 Checklist Reproducibility (dievaluasi otomatis dari hasil run)}\\
\toprule
Item & Pass\\
\midrule
Seed eksplisit tercatat & TRUE\\
Checksum SHA-256 data mentah tercatat & TRUE\\
Preprocessing scripted (make\_frames) & TRUE\\
Standardisasi dihitung di dalam fold (tanpa kebocoran) & TRUE\\
Prior eksplisit + prior predictive check & TRUE\\
\addlinespace
Komputasi divalidasi (analitik vs Gibbs vs HMC, < 3 MCSE) & TRUE\\
Likelihood manual = likelihood Stan & TRUE\\
Semua model: R-hat <= 1.01, ESS >= 400, 0 divergence & TRUE\\
Stabil terhadap seed (semua \&\#124;z\&\#124; < 3 MCSE) & TRUE\\
Stabil terhadap urutan data (semua \&\#124;z\&\#124; < 3 MCSE) & TRUE\\
\addlinespace
Kesimpulan prediktif stabil antar 5 seed (rentang P(R>20 \&\#124; Fr=0.40) < 0.05) & TRUE\\
Validasi per-hull dipakai untuk klaim prediktif & TRUE\\
Objek fit tersimpan (fits/*.rds) & TRUE\\
\bottomrule
\end{longtable}
\endgroup{}

\clearpage

\section{BAB 5. KESIMPULAN DAN SARAN}\label{bab-5.-kesimpulan-dan-saran}

\subsection{5.1 Kesimpulan}\label{kesimpulan}

\begin{enumerate}
\def\labelenumi{\arabic{enumi}.}
\tightlist
\item
  \textbf{Komputasi terverifikasi dan stabil secara substantif.}
  Posterior MCMC cocok dengan solusi analitik dan Gibbs (≤ 1.02 MCSE).
  Kesimpulan prediktif tidak berubah antar lima seed maupun urutan data,
  dan prior yang masih masuk akal tidak mengubah kesimpulan. Namun MCSE
  dari satu run meremehkan variabilitas antar-run untuk SD efek hull
  (≈0.5×), sehingga audit reproducibility perlu memakai beberapa seed.
\item
  \textbf{Keputusan kualitas data dan preprocessing mendominasi hasil.}
  Baseline Gaussian pada skala mentah konvergen sempurna tetapi salah
  spesifikasi. Transformasi log dan bentuk fungsi Froude yang fleksibel
  memperbaiki ELPD dan RMSE secara drastis. Prior yang salah skala
  mengubah kesimpulan tanpa peringatan konvergensi, sehingga prior
  predictive check wajib dilakukan.
\item
  \textbf{Skema validasi adalah sumber optimisme terbesar} dan bahkan
  mengubah model yang terpilih. Untuk estimand utama (prediksi lambung
  baru), validasi harus dilakukan per hull.
\item
  \textbf{Informatif untuk prediksi, terbatas untuk inferensi geometri.}
  Kurva resistance--Froude, heteroskedastisitas, dan variasi antar-hull
  diestimasi dengan presisi tinggi. Interval prediksi 95\% M3
  terkalibrasi baik untuk hull baru (cakupan 0.94), sedangkan M6 sedikit
  konservatif (0.99). Namun arah efek tiap rasio geometri tidak dapat
  dipastikan (22 hull, VIF tinggi).
\end{enumerate}

\subsection{5.2 Keterbatasan}\label{keterbatasan}

\begin{enumerate}
\def\labelenumi{\arabic{enumi}.}
\tightlist
\item
  Hanya 22 hull yang semuanya turunan satu \emph{parent hull};
  generalisasi ke keluarga desain lain tidak didukung.
\item
  Geometri divariasikan satu per satu di sekitar \emph{parent hull},
  sehingga efeknya lokal dan interaksi geometri tidak teridentifikasi.
\item
  Froude hanya 14 level (0.125--0.45); prediksi di luar rentang ini
  adalah ekstrapolasi.
\item
  Transformasi \(\log(1+y)\) masih mengizinkan nilai −1 \textless{} y
  \textless{} 0.
\item
  Galat pengukuran \emph{towing tank} tidak terdokumentasi.
\item
  Efek geometri adalah asosiasi dalam desain eksperimen ini, bukan klaim
  kausal umum.
\end{enumerate}

\subsection{5.3 Saran dan Pengembangan}\label{saran-dan-pengembangan}

\begin{enumerate}
\def\labelenumi{\arabic{enumi}.}
\tightlist
\item
  Gunakan validasi per hull (atau per kelompok) untuk setiap klaim
  prediktif pada data berkelompok, dan laporkan optimisme skema acak
  sebagai bagian audit.
\item
  Uji likelihood yang hanya berlaku pada \(y>0\) (lognormal atau Gamma)
  untuk menghilangkan sisa pelanggaran support.
\item
  Untuk prediksi desain baru, gunakan model berbasis spline Froude yang
  lebih sederhana; gunakan M6 untuk menjelaskan struktur data.
\item
  Tambahkan data dari keluarga lambung lain agar efek geometri dapat
  diidentifikasi.
\end{enumerate}

\newpage

\section{DAFTAR PUSTAKA}\label{daftar-pustaka}

\begin{enumerate}
\def\labelenumi{\arabic{enumi}.}
\tightlist
\item
  \textbf{Bürkner, P.-C. (2017).} brms: An R package for Bayesian
  multilevel models using Stan. \emph{Journal of Statistical Software,
  80}(1), 1--28. \url{https://doi.org/10.18637/jss.v080.i01}
\item
  \textbf{Gelman, A., Carlin, J. B., Stern, H. S., Dunson, D. B.,
  Vehtari, A., \& Rubin, D. B. (2013).} \emph{Bayesian data analysis}
  (3rd ed.). CRC Press.
\item
  \textbf{Gelman, A., Vehtari, A., Simpson, D., Margossian, C. C.,
  Carpenter, B., Yao, Y., Kennedy, L., Gabry, J., Bürkner, P.-C., \&
  Modrák, M. (2020).} Bayesian workflow. \emph{arXiv}.
  \url{https://arxiv.org/abs/2011.01808}
\item
  \textbf{Gerritsma, J., Onnink, R., \& Versluis, A. (1981).} Geometry,
  resistance and stability of the Delft systematic yacht hull series.
  \emph{International Shipbuilding Progress, 28}(328), 276--297.
\item
  \textbf{Hoff, P. D. (2009).} \emph{A first course in Bayesian
  statistical methods}. Springer.
  \url{https://doi.org/10.1007/978-0-387-92407-6}
\item
  \textbf{Roberts, D. R., Bahn, V., Ciuti, S., Boyce, M. S., Elith, J.,
  Guillera-Arroita, G., \ldots{} \& Dormann, C. F. (2017).}
  Cross-validation strategies for data with temporal, spatial,
  hierarchical, or phylogenetic structure. \emph{Ecography, 40}(8),
  913--929. \url{https://doi.org/10.1111/ecog.02881}
\item
  \textbf{Vehtari, A., Gelman, A., \& Gabry, J. (2017).} Practical
  Bayesian model evaluation using leave-one-out cross-validation and
  WAIC. \emph{Statistics and Computing, 27}, 1413--1432.
  \url{https://doi.org/10.1007/s11222-016-9696-4}
\item
  \textbf{Vehtari, A., Gelman, A., Simpson, D., Carpenter, B., \&
  Bürkner, P.-C. (2021).} Rank-normalization, folding, and localization:
  An improved \(\hat R\) for assessing convergence of MCMC.
  \emph{Bayesian Analysis, 16}(2), 667--718.
  \url{https://doi.org/10.1214/20-BA1221}
\item
  \textbf{Gerritsma, J., Onnink, R., \& Versluis, A. (1981).}
  \emph{Yacht Hydrodynamics} {[}Dataset{]}. UCI Machine Learning
  Repository. \url{https://doi.org/10.24432/C5XG7R}
\end{enumerate}

\section{ACKNOWLEDGEMENT}\label{acknowledgement}

Sebagai wujud komitmen terhadap integritas akademik, kami menyatakan
bahwa dalam penyelesaian project ini kami memanfaatkan \emph{Artificial
Intelligence}, yaitu Claude (Anthropic).

Penggunaan AI dibatasi sebagai alat pendukung (\emph{assistive tool}):
bantuan teknis dalam penulisan dan pemeriksaan kode R/Python, penyusunan
struktur dan keterbacaan laporan, serta mitra diskusi dalam memvalidasi
pemahaman terhadap \emph{Bayesian workflow}. Pemeriksaan kode dengan
bantuan AI juga menemukan kesalahan skala prior pada versi awal analisis
(Subbab 4.2), yang kemudian diperbaiki dan didokumentasikan sebagai
temuan audit.

Seluruh angka dalam laporan ini dihasilkan oleh kode yang dapat
dijalankan ulang tanpa AI (\texttt{analisis\_yacht\_final.R}). Tidak ada
angka yang diambil dari teks keluaran AI. Kami menegaskan bahwa
verifikasi data, penalaran analitis, dan penarikan kesimpulan akhir
dilakukan secara mandiri dan sepenuhnya menjadi tanggung jawab kami
sebagai penulis.

\section{LAMPIRAN}\label{lampiran}

\subsubsection{Lampiran A: Minimum Reporting Checklist (Bayesian
Regression)}\label{lampiran-a-minimum-reporting-checklist-bayesian-regression}

\begingroup\fontsize{8.5}{10.5}\selectfont

\begin{longtable}[t]{ll}
\caption{\label{tab:lamp_report}Tabel A.1 Minimum Reporting Checklist}\\
\toprule
Item & Isi\\
\midrule
Likelihood \& link & Gaussian pada log(1+Resistance) terstandardisasi; geometri linear + B-spline Froude (df = 6) + efek acak hull; log-link untuk sigma(Froude)\\
Prior & Tabel 3.3 (divalidasi dengan prior predictive check, Subbab 4.2)\\
Software & R version 4.5.2 (2025-10-31 ucrt); brms 2.23.0; rstan 2.32.7; loo 2.10.1\\
Chains / iter / warm-up & 4 / 2000 / 1000; adapt\_delta 0.95; seed 2026\\
Diagnostik M6 & max R-hat 1.004; min ESS bulk 1,036; min ESS tail 1,644; divergences 0\\
\addlinespace
Posterior summary & Tabel 4.10: median, mean, mode, SD, 95\% ETI \& HDI, posterior probability\\
PPC & Gambar 4.3 \& 4.8; Tabel 4.5 \& 4.9\\
Sensitivitas & Subbab 4.9: prior x0.5/x1/x2.5; Gaussian vs Student-t; homo vs hetero; independen vs hierarkis\\
Perbandingan prediktif & PSIS-LOO \& WAIC (Jacobian); leave-hull-out CV M6: RMSE 2.09, coverage 0.99\\
\bottomrule
\end{longtable}
\endgroup{}

\subsubsection{Lampiran B: Cara
Mereproduksi}\label{lampiran-b-cara-mereproduksi}

\begin{enumerate}
\def\labelenumi{\arabic{enumi}.}
\tightlist
\item
  Letakkan \texttt{yacht\_hydrodynamics.data},
  \texttt{analisis\_yacht\_final.R}, dan file \texttt{.Rmd} ini dalam
  satu folder.
\item
  Pasang paket:
  \texttt{install.packages(c("brms",\ "rstan",\ "loo",\ "posterior",\ "dplyr",\ "tidyr",\ "kableExtra",\ "patchwork",\ "bayesplot",\ "digest",\ "jsonlite"))}.
\item
  \emph{Knit} file \texttt{.Rmd}. Run pertama menghitung semua model
  (±40--60 menit); hasil disimpan di \texttt{fits/} dan \texttt{cache/}
  sehingga knit berikutnya jauh lebih cepat dan menghasilkan angka
  identik. Hapus kedua folder tersebut untuk menghitung ulang dari nol.
\item
  Checksum SHA-256 data mentah:
  00dfecc0fc01ddd4c90b558a3ac11b246df8ebcfea130724223475a9a67f0ea1.
  Manifest lengkap tersimpan di
  \texttt{reproducibility\_manifest\_R.json}.
\end{enumerate}

\subsubsection{Lampiran C: Informasi Sesi
R}\label{lampiran-c-informasi-sesi-r}

\begin{verbatim}
R version 4.5.2 (2025-10-31 ucrt)
Platform: x86_64-w64-mingw32/x64
Running under: Windows 11 x64 (build 26200)

Matrix products: default
  LAPACK version 3.12.1

locale:
[1] LC_COLLATE=English_Indonesia.utf8  LC_CTYPE=English_Indonesia.utf8   
[3] LC_MONETARY=English_Indonesia.utf8 LC_NUMERIC=C                      
[5] LC_TIME=English_Indonesia.utf8    

time zone: Asia/Jakarta
tzcode source: internal

attached base packages:
[1] splines   stats     graphics  grDevices utils     datasets  methods  
[8] base     

other attached packages:
 [1] tidyr_1.3.1        dplyr_1.1.4        posterior_1.7.0    loo_2.10.1        
 [5] rstan_2.32.7       StanHeaders_2.32.7 brms_2.23.0        Rcpp_1.1.0        
 [9] bayesplot_1.16.0   patchwork_1.3.2    ggplot2_4.0.1      kableExtra_1.4.0  
[13] knitr_1.50        

loaded via a namespace (and not attached):
 [1] gtable_0.3.6         tensorA_0.36.2.1     QuickJSR_1.11.0     
 [4] xfun_0.55            inline_0.3.21        lattice_0.22-7      
 [7] vctrs_0.6.5          tools_4.5.2          generics_0.1.4      
[10] stats4_4.5.2         parallel_4.5.2       sandwich_3.1-1      
[13] tibble_3.3.0         pkgconfig_2.0.3      Matrix_1.7-4        
[16] checkmate_2.3.4      RColorBrewer_1.1-3   S7_0.2.1            
[19] distributional_0.9.0 RcppParallel_6.2.1   lifecycle_1.0.4     
[22] compiler_4.5.2       farver_2.1.2         stringr_1.6.0       
[25] textshaping_1.0.4    Brobdingnag_1.3-1    codetools_0.2-20    
[28] htmltools_0.5.9      yaml_2.3.12          pillar_1.11.1       
[31] MASS_7.3-65          bridgesampling_1.2-1 abind_1.4-8         
[34] multcomp_1.4-29      nlme_3.1-168         tidyselect_1.2.1    
[37] digest_0.6.39        mvtnorm_1.3-3        stringi_1.8.7       
[40] reshape2_1.4.5       purrr_1.2.0          labeling_0.4.3      
[43] fastmap_1.2.0        grid_4.5.2           cli_3.6.5           
[46] magrittr_2.0.4       pkgbuild_1.4.8       survival_3.8-3      
[49] TH.data_1.1-5        withr_3.0.2          scales_1.4.0        
[52] backports_1.5.0      estimability_1.5.1   rmarkdown_2.30      
[55] matrixStats_1.5.0    emmeans_2.0.2        gridExtra_2.3       
[58] zoo_1.8-15           coda_0.19-4.1        evaluate_1.0.5      
[61] viridisLite_0.4.2    rstantools_2.7.1     rlang_1.1.7         
[64] xtable_1.8-4         glue_1.8.0           xml2_1.5.1          
[67] jsonlite_2.0.0       svglite_2.2.2        rstudioapi_0.17.1   
[70] plyr_1.8.9           R6_2.6.1             systemfonts_1.3.1   
\end{verbatim}

\subsubsection{Lampiran D: Source Code Analisis Utama
(R)}\label{lampiran-d-source-code-analisis-utama-r}

Berikut adalah kode R lengkap yang digunakan untuk seluruh analisis
statistik dan pemodelan:

\begin{verbatim}
# =============================================================================
# analisis_yacht.R
# Regresi Bayesian + audit reproducibility & data quality - Yacht Hydrodynamics
# Seluruh pipeline dapat dijalankan ulang tanpa AI:  source("analisis_yacht.R")
# Hasil fit disimpan di folder fits/ dan cache/ agar knit ulang cepat & identik.
# =============================================================================

suppressPackageStartupMessages({
  library(brms)      # regresi Bayesian (Stan/HMC)
  library(rstan)     # model Stan konjugat untuk validasi komputasi
  library(loo)       # PSIS-LOO dan WAIC
  library(posterior) # ringkasan draws, R-hat, ESS, MCSE
  library(dplyr)
  library(tidyr)
  library(splines)
})

SEED <- 2026
set.seed(SEED)
N_CORES <- min(4, parallel::detectCores())
options(mc.cores = N_CORES)
rstan_options(auto_write = TRUE)
dir.create("fits", showWarnings = FALSE)
dir.create("cache", showWarnings = FALSE)

# Hasil mahal disimpan sebagai .rds; hapus folder cache/ dan fits/ untuk menghitung ulang dari nol
cache_rds <- function(name, expr) {
  f <- file.path("cache", paste0(name, ".rds"))
  if (file.exists(f)) return(readRDS(f))
  val <- expr
  saveRDS(val, f)
  val
}

SAMPLER    <- list(chains = 4, iter = 2000, warmup = 1000)   # 4 x 1000 draws
CV_SAMPLER <- list(chains = 2, iter = 1000, warmup = 500)    # konfigurasi ringan khusus CV
CONTROL    <- list(adapt_delta = 0.95)

# -----------------------------------------------------------------------------
# 1. DATA: load, nama kolom sesuai dokumentasi UCI 243, checksum
# -----------------------------------------------------------------------------
COLS <- c("LCB", "Prismatic_Coef", "LD_Ratio", "BD_Ratio", "LB_Ratio", "Froude_No", "Resistance")
PRED <- COLS[1:6]
GEOM <- COLS[1:5]
RAW_PATH <- "yacht_hydrodynamics.data"
df <- read.table(RAW_PATH, col.names = COLS)
RAW_SHA256 <- digest::digest(file = RAW_PATH, algo = "sha256")

geom_key <- do.call(paste, df[GEOM])
df$hull_id <- as.integer(factor(geom_key, levels = unique(geom_key)))

sdp  <- function(x) sqrt(mean((x - mean(x))^2))                   # SD populasi (ddof = 0)
skew <- function(x) mean((x - mean(x))^3) / sdp(x)^3

# -----------------------------------------------------------------------------
# 2. AUDIT KUALITAS DATA
# -----------------------------------------------------------------------------
rows_per_hull <- table(df$hull_id)
iqr_out <- sapply(COLS, function(v) {
  q <- quantile(df[[v]], c(.25, .75)); w <- diff(q)
  sum(df[[v]] < q[1] - 1.5 * w | df[[v]] > q[2] + 1.5 * w)
})
AUDIT <- list(
  n_rows = nrow(df), n_cols = length(COLS), n_missing = sum(is.na(df[COLS])),
  n_dup = sum(duplicated(df[COLS])), n_hull = length(unique(df$hull_id)),
  rows_per_hull = unique(as.integer(rows_per_hull)), n_froude = length(unique(df$Froude_No)),
  skew_raw = skew(df$Resistance), skew_log = skew(log1p(df$Resistance)),
  shapiro_p = shapiro.test(df$Resistance)$p.value, iqr_out = iqr_out,
  range = t(sapply(COLS, function(v) c(min = min(df[[v]]), max = max(df[[v]]),
                                       mean = mean(df[[v]]), sd = sd(df[[v]]))))
)

# -----------------------------------------------------------------------------
# 3. PREPROCESSING EKSPLISIT (statistik dihitung dari data TRAIN saja)
# -----------------------------------------------------------------------------
make_frames <- function(train, test = NULL) {
  sc <- lapply(PRED, function(v) c(mean(train[[v]]), sdp(train[[v]])))
  names(sc) <- PRED
  sc$y  <- c(mean(train$Resistance), sdp(train$Resistance))
  ly    <- log1p(train$Resistance)
  sc$ly <- c(mean(ly), sdp(ly))
  basis <- bs((train$Froude_No - sc$Froude_No[1]) / sc$Froude_No[2], df = 6)  # B-spline kubik, knot dari train
  tf <- function(d) {
    o <- as.data.frame(lapply(setNames(PRED, paste0("z_", PRED)),
                              function(v) (d[[v]] - sc[[v]][1]) / sc[[v]][2]))
    o$zF2 <- o$z_Froude_No^2
    o$zF3 <- o$z_Froude_No^3
    B <- predict(basis, o$z_Froude_No)
    colnames(B) <- paste0("F", 1:6)
    o <- cbind(o, as.data.frame(B))
    o$y_z  <- (d$Resistance - sc$y[1]) / sc$y[2]
    o$ly_z <- (log1p(d$Resistance) - sc$ly[1]) / sc$ly[2]
    o$hull <- sprintf("H%02d", d$hull_id)
    o$hull_id <- d$hull_id
    o$Froude_No <- d$Froude_No
    o$Resistance <- d$Resistance
    o
  }
  list(train = tf(train), test = if (!is.null(test)) tf(test), sc = sc, tf = tf)
}

full <- df[c(PRED, "Resistance", "hull_id")]
FR <- make_frames(full)
mdl <- FR$train
SCALE <- FR$sc

to_raw <- function(z, resp, sc) if (resp == "raw") z * sc$y[2] + sc$y[1] else expm1(z * sc$ly[2] + sc$ly[1])
jacobian <- function(frame, resp, sc) {
  if (resp == "raw") rep(-log(sc$y[2]), nrow(frame)) else -log(sc$ly[2]) - log1p(frame$Resistance)
}

# -----------------------------------------------------------------------------
# 4. SPESIFIKASI MODEL & PRIOR
# -----------------------------------------------------------------------------
GEO <- paste0("z_", GEOM, collapse = " + ")
LIN <- paste(GEO, "+ z_Froude_No")
SPL <- paste0("F", 1:6, collapse = " + ")

SPECS <- list(
  M0 = list(label = "M0 Baseline: Gaussian, raw target (Variant 1)", f = paste("y_z ~", LIN),
            resp = "raw", fam = "gaussian", froude = "z_Froude_No", spline = FALSE, hetero = FALSE, hier = FALSE),
  M1 = list(label = "M1 Log target, linear (Variant 2)", f = paste("ly_z ~", LIN),
            resp = "log", fam = "gaussian", froude = "z_Froude_No", spline = FALSE, hetero = FALSE, hier = FALSE),
  M2 = list(label = "M2 Log target + Froude^2, ^3 (Variant 3)", f = paste("ly_z ~", LIN, "+ zF2 + zF3"),
            resp = "log", fam = "gaussian", froude = c("z_Froude_No", "zF2", "zF3"), spline = FALSE, hetero = FALSE, hier = FALSE),
  M3 = list(label = "M3 Log target + Froude spline", f = paste("ly_z ~", GEO, "+", SPL),
            resp = "log", fam = "gaussian", froude = character(0), spline = TRUE, hetero = FALSE, hier = FALSE),
  M4 = list(label = "M4 M3 + heteroskedastic sigma(Froude)", f = paste("ly_z ~", GEO, "+", SPL),
            resp = "log", fam = "gaussian", froude = character(0), spline = TRUE, hetero = TRUE, hier = FALSE),
  M5 = list(label = "M5 M3 with Student-t likelihood", f = paste("ly_z ~", GEO, "+", SPL),
            resp = "log", fam = "student", froude = character(0), spline = TRUE, hetero = FALSE, hier = FALSE),
  M6 = list(label = "M6 M4 + hull random intercept (final)", f = paste("ly_z ~", GEO, "+", SPL, "+ (1 | hull)"),
            resp = "log", fam = "gaussian", froude = character(0), spline = TRUE, hetero = TRUE, hier = TRUE)
)

make_formula <- function(sp) {
  if (sp$hetero) bf(as.formula(sp$f), sigma ~ z_Froude_No) else bf(as.formula(sp$f))
}

# s = pengali skala SD prior (1 = dasar, 0.5 lebih ketat, 2.5 lebih lebar) untuk analisis sensitivitas
make_prior <- function(sp, s = 1) {
  p <- prior_string(sprintf("normal(0, %g)", 1 * s), class = "Intercept") +
       prior_string(sprintf("normal(0, %g)", 0.5 * s), class = "b")                  # geometri
  for (v in sp$froude) p <- p + prior_string(sprintf("normal(0, %g)", 1 * s), class = "b", coef = v)
  if (sp$spline) for (k in 1:6) p <- p + prior_string(sprintf("normal(0, %g)", 1.5 * s), class = "b", coef = paste0("F", k))
  if (sp$hetero) {
    p <- p + prior_string(sprintf("normal(-1.5, %g)", 0.75 * s), class = "Intercept", dpar = "sigma") +
             prior_string(sprintf("normal(0, %g)", 0.5 * s), class = "b", dpar = "sigma")
  } else {
    p <- p + prior_string(sprintf("exponential(%g)", 1 / s), class = "sigma")          # rate = 1/s
  }
  if (sp$fam == "student") p <- p + prior(gamma(2, 0.1), class = "nu")               # shape-rate
  if (sp$hier) p <- p + prior_string(sprintf("normal(0, %g)", 0.25 * s), class = "sd")  # half-normal
  p
}

fit_model <- function(name, data = mdl, s = 1, seed = SEED, tag = name, sampler = SAMPLER) {
  sp <- SPECS[[name]]
  invisible(gc())                                   # bebaskan memori sebelum kompilasi Stan
  brm(make_formula(sp), data = data,
      family = if (sp$fam == "student") student() else gaussian(),
      prior = make_prior(sp, s), chains = sampler$chains, iter = sampler$iter,
      warmup = sampler$warmup, cores = N_CORES, seed = seed, control = CONTROL,
      refresh = 0, silent = 2, file = file.path("fits", tag), file_refit = "on_change")
}

MAIN_PAR <- function(fit) grep("^b_|^sigma$|^sd_|^nu$", variables(fit), value = TRUE)

diag_row <- function(fit, name) {
  t0 <- sum(rstan::get_elapsed_time(fit$fit))
  s <- summarise_draws(as_draws_array(fit, variable = MAIN_PAR(fit)),
                       "mean", "sd", "rhat", "ess_bulk", "ess_tail", "mcse_mean")
  np <- nuts_params(fit)
  data.frame(model = name, sampling_s = round(t0, 1), max_rhat = max(s$rhat),
             min_ess_bulk = min(s$ess_bulk), min_ess_tail = min(s$ess_tail),
             max_mcse_over_sd = max(s$mcse_mean / s$sd),
             divergences = sum(np$Value[np$Parameter == "divergent__"]))
}

loo_raw <- function(fit, name, data = mdl) {
  ll <- log_lik(fit)
  ll <- sweep(ll, 2, jacobian(data, SPECS[[name]]$resp, SCALE), "+")   # pindah ke satuan asli Resistance
  chain_id <- rep(seq_len(nchains(fit)), each = ndraws(fit) / nchains(fit))
  list(loo = loo(ll, r_eff = relative_eff(exp(ll), chain_id = chain_id)), waic = waic(ll))
}

# -----------------------------------------------------------------------------
# 5. PRIOR PREDICTIVE CHECK (simulasi eksplisit dari prior, materi Bag. 5.3)
# -----------------------------------------------------------------------------
prior_predictive <- function(name, S = 1000, wide = FALSE) {
  sp <- SPECS[[name]]
  set.seed(SEED)
  Xg <- scale(as.matrix(mdl[paste0("z_", GEOM)]), scale = FALSE)          # brms: intercept pada prediktor terpusat
  zF <- mdl$z_Froude_No
  out <- matrix(NA_real_, S, nrow(mdl))
  for (s in 1:S) {
    ic  <- rnorm(1, 0, if (wide) 2.5 else 1)
    mu  <- ic + Xg %*% rnorm(length(GEOM), 0, if (wide) 1 else 0.5)
    if (sp$spline) {
      Fm <- scale(as.matrix(mdl[paste0("F", 1:6)]), scale = FALSE)
      mu <- mu + Fm %*% rnorm(6, 0, if (wide) 2.5 else 1.5)
    } else {
      mu <- mu + (zF - mean(zF)) * rnorm(1, 0, 1)
    }
    if (sp$hier) {
      u <- rnorm(AUDIT$n_hull, 0, abs(rnorm(1, 0, if (wide) 1 else 0.25)))
      mu <- mu + u[mdl$hull_id]
    }
    sig <- if (sp$hetero) exp(rnorm(1, if (wide) 0 else -1.5, if (wide) 2 else 0.75) +
                              rnorm(1, 0, if (wide) 1 else 0.5) * zF) else rexp(1, 1)
    out[s, ] <- rnorm(nrow(mdl), mu, sig)
  }
  to_raw(out, sp$resp, SCALE)
}
PRIOR_PRED <- cache_rds("prior_predictive", list(
  M0 = prior_predictive("M0"), M6 = prior_predictive("M6"), M6_initial = prior_predictive("M6", wide = TRUE)))

pp_summary <- function(v) {
  v <- v[is.finite(v)]
  c(share_negative = mean(v < 0), share_below_minus_0.5 = mean(v < -0.5), share_gt_1000 = mean(v > 1000),
    q01 = unname(quantile(v, .01)), q50 = median(v), q95 = unname(quantile(v, .95)), q99 = unname(quantile(v, .99)))
}

# -----------------------------------------------------------------------------
# 6. PRIOR SALAH SKALA (reproduksi kesalahan versi awal, materi Bag. 7/10A)
# -----------------------------------------------------------------------------
fit_old <- brm(as.formula(paste("Resistance ~", paste(PRED, collapse = " + "))), data = df,
               prior = prior_string(sprintf("normal(%g, %g)", mean(df$Resistance), 2.5 * sd(df$Resistance)), class = "Intercept") +
                       prior(normal(0, 1), class = "b") +
                       prior_string(sprintf("normal(0, %g)", sd(df$Resistance)), class = "sigma"),
               chains = CV_SAMPLER$chains, iter = CV_SAMPLER$iter, warmup = CV_SAMPLER$warmup, cores = N_CORES,
               seed = SEED, refresh = 0, silent = 2, file = "fits/old_misscaled_prior", file_refit = "on_change")
Xraw <- cbind(1, as.matrix(df[PRED]))
b_ols_raw <- drop(solve(crossprod(Xraw), crossprod(Xraw, df$Resistance)))
OLD <- list(ols = b_ols_raw[["Froude_No"]], old = fixef(fit_old)["Froude_No", ],
            prior_span95 = 1.96 * diff(range(df$Froude_No)))

# -----------------------------------------------------------------------------
# 7. VALIDASI KOMPUTASI: analitik konjugat vs Gibbs manual vs HMC (Stan)
# -----------------------------------------------------------------------------
X <- cbind(Intercept = 1, as.matrix(mdl[paste0("z_", PRED)]))
y <- mdl$y_z
n <- nrow(X); p <- ncol(X)
m0 <- rep(0, p); V0 <- diag(p); a0 <- 2; b0 <- 1
V0i <- solve(V0)
Vn <- solve(V0i + crossprod(X))
mn <- drop(Vn %*% (V0i %*% m0 + crossprod(X, y)))
an <- a0 + n / 2
bn <- b0 + 0.5 * drop(crossprod(y) + t(m0) %*% V0i %*% m0 - t(mn) %*% solve(Vn) %*% mn)
sd_analytic <- sqrt(diag(Vn) * bn / (an - 1))
b_ols <- drop(solve(crossprod(X), crossprod(X, y)))

GIBBS <- cache_rds("gibbs", {
  set.seed(SEED)
  B <- 12000; burn <- 2000
  beta <- matrix(0, B, p); s2 <- rep(1, B)
  L <- t(chol(Vn))
  for (it in 2:B) {
    beta[it, ] <- mn + sqrt(s2[it - 1]) * drop(L %*% rnorm(p))
    r <- y - X %*% beta[it, ]
    dp <- beta[it, ] - m0
    s2[it] <- 1 / rgamma(1, shape = a0 + (n + p) / 2,
                         rate = b0 + 0.5 * drop(crossprod(r) + t(dp) %*% V0i %*% dp))  # 1/sigma2 ~ Gamma(shape, rate)
  }
  list(beta = beta[-(1:burn), ], s2 = s2[-(1:burn)])
})

NIG_STAN <- "
data { int<lower=1> N; int<lower=1> P; matrix[N, P] X; vector[N] y; real a0; real b0; }
parameters { real<lower=0> sigma2; vector[P] beta_raw; }
transformed parameters { vector[P] beta = sqrt(sigma2) * beta_raw; }
model { sigma2 ~ inv_gamma(a0, b0); beta_raw ~ std_normal(); y ~ normal(X * beta, sqrt(sigma2)); }
"
HMC_NIG <- cache_rds("hmc_nig", {
  sm <- stan_model(model_code = NIG_STAN)
  fit <- sampling(sm, data = list(N = n, P = p, X = X, y = y, a0 = a0, b0 = b0),
                  chains = 4, iter = 3000, warmup = 1000, seed = SEED, refresh = 0)
  summarise_draws(as_draws_array(fit), "mean", "sd", "mcse_mean", "rhat", "ess_bulk")
})
hmc_beta <- HMC_NIG[grepl("^beta\\[", HMC_NIG$variable), ]
COMP <- data.frame(
  OLS = b_ols, Analytic_mean = mn, Analytic_SD = sd_analytic, Gibbs_mean = colMeans(GIBBS$beta),
  HMC_mean = hmc_beta$mean, HMC_MCSE = hmc_beta$mcse_mean, row.names = colnames(X))
COMP$HMC_minus_Analytic_in_MCSE <- (COMP$HMC_mean - COMP$Analytic_mean) / COMP$HMC_MCSE
SIGMA2 <- c(Analytic = bn / (an - 1), Gibbs = mean(GIBBS$s2),
            HMC = HMC_NIG$mean[HMC_NIG$variable == "sigma2"])

# log-likelihood manual (materi Bag. 4) untuk diverifikasi terhadap Stan
loglik_gaussian <- function(beta, sigma, X, y) sum(dnorm(y, drop(X %*% beta), sigma, log = TRUE))

# -----------------------------------------------------------------------------
# 8. FIT MODEL M0 - M6
# -----------------------------------------------------------------------------
FITS <- lapply(setNames(names(SPECS), names(SPECS)), fit_model)
DIAG <- bind_rows(lapply(names(FITS), function(m) diag_row(FITS[[m]], m)))
DIAG$verdict <- ifelse(DIAG$max_rhat <= 1.01 & DIAG$min_ess_bulk >= 400 & DIAG$min_ess_tail >= 400 &
                         DIAG$divergences == 0, "OK", "CHECK")

# verifikasi likelihood manual pada satu posterior draw M0
d0 <- as_draws_df(FITS$M0)
LL_CHECK <- c(
  manual = loglik_gaussian(unlist(d0[1, paste0("b_", c("Intercept", paste0("z_", PRED)))]), d0$sigma[1], X, y) +
    n * (-log(SCALE$y[2])),
  stan = sum(log_lik(FITS$M0, draw_ids = 1)) + n * (-log(SCALE$y[2])))

# -----------------------------------------------------------------------------
# 9. LOO / WAIC (koreksi Jacobian ke satuan asli Resistance)
# -----------------------------------------------------------------------------
# Kompatibilitas cache loo:
# Object loo yang tersimpan dari versi package lama dapat memiliki struktur
# diagnostik Pareto-k yang berbeda. Ambil Pareto-k dari struktur yang tersedia,
# dan bila cache lama tidak menyimpannya, hitung ulang LOOS dari fit saat ini.
extract_pareto_k <- function(x) {
  if (is.null(x)) return(NULL)

  # Struktur utama pada loo modern
  if (!is.null(x$diagnostics) && !is.null(x$diagnostics$pareto_k)) {
    return(as.numeric(x$diagnostics$pareto_k))
  }

  # Fallback yang tetap tersedia pada objek loo dengan pointwise diagnostics
  if (!is.null(x$pointwise) &&
      "influence_pareto_k" %in% colnames(x$pointwise)) {
    return(as.numeric(x$pointwise[, "influence_pareto_k"]))
  }

  # Fallback terakhir bila psis object tersimpan
  if (!is.null(x$psis_object)) {
    out <- tryCatch(loo::pareto_k_values(x$psis_object),
                    error = function(e) NULL)
    if (!is.null(out)) return(as.numeric(out))
  }

  NULL
}

build_loos <- function() {
  lapply(setNames(names(FITS), names(FITS)),
         function(m) loo_raw(FITS[[m]], m))
}

LOOS <- cache_rds("loos", build_loos())

# Cache loo lama mungkin sah untuk ELPD/WAIC tetapi tidak menyimpan
# diagnostik Pareto-k yang dikenali oleh versi loo yang sedang dipakai.
# Dalam kondisi ini, cache khusus LOO dibangun ulang dari posterior yang ada.
has_pareto_k <- vapply(LOOS, function(z) {
  pk <- extract_pareto_k(z$loo)
  length(pk) > 0 && any(is.finite(pk))
}, logical(1))

if (!all(has_pareto_k)) {
  unlink(file.path("cache", "loos.rds"))
  LOOS <- cache_rds("loos", build_loos())
}

loo_cmp <- loo_compare(lapply(LOOS, `[[`, "loo"))
LOO_TAB <- data.frame(model = rownames(loo_cmp), loo_cmp[, c("elpd_loo", "se_elpd_loo", "p_loo", "elpd_diff", "se_diff")],
                      row.names = NULL)
LOO_TAB$elpd_waic <- sapply(LOO_TAB$model, function(m) LOOS[[m]]$waic$estimates["elpd_waic", "Estimate"])

LOO_TAB$n_pareto_k_gt_0.7 <- sapply(LOO_TAB$model, function(m) {
  pk <- extract_pareto_k(LOOS[[m]]$loo)
  if (is.null(pk)) return(NA_integer_)
  sum(pk > 0.7, na.rm = TRUE)
})

LOO_TAB$diff_over_se <- ifelse(LOO_TAB$se_diff > 0, LOO_TAB$elpd_diff / LOO_TAB$se_diff, NA)

# -----------------------------------------------------------------------------
# 10. PPC & RESIDUAL (satuan asli)
# -----------------------------------------------------------------------------
ppc_raw <- function(name, ndraws = 1000) {
  set.seed(SEED)
  to_raw(posterior_predict(FITS[[name]], ndraws = ndraws), SPECS[[name]]$resp, SCALE)
}
STATS <- list(mean = mean, sd = sd, max = max, skewness = skew, `share < 0` = function(v) mean(v < 0))
ppc_table <- function(yrep, yobs) {
  do.call(rbind, lapply(names(STATS), function(k) {
    rep <- apply(yrep, 1, STATS[[k]]); obs <- STATS[[k]](yobs)
    data.frame(statistic = k, observed = obs, rep_median = median(rep), rep_2.5 = unname(quantile(rep, .025)),
               rep_97.5 = unname(quantile(rep, .975)), P_rep_ge_obs = mean(rep >= obs))
  }))
}
YOBS <- mdl$Resistance
YREP <- cache_rds("yrep", list(M0 = ppc_raw("M0"), M6 = ppc_raw("M6")))
PPC <- list(M0 = ppc_table(YREP$M0, YOBS), M6 = ppc_table(YREP$M6, YOBS))

MU_M0 <- colMeans(to_raw(posterior_linpred(FITS$M0), "raw", SCALE))
MU6 <- colMeans(posterior_linpred(FITS$M6))
SIG6 <- colMeans(posterior_linpred(FITS$M6, dpar = "sigma", transform = TRUE))
RES6 <- (mdl$ly_z - MU6) / SIG6

# VIF dan korelasi posterior (materi Bag. 16, 24.4)
vif_table <- function(M) sapply(seq_len(ncol(M)), function(i) 1 / (1 - summary(lm(M[, i] ~ M[, -i]))$r.squared))
VIF <- list(geometry = setNames(vif_table(as.matrix(mdl[paste0("z_", PRED)])), PRED),
            poly = setNames(vif_table(as.matrix(mdl[c("z_Froude_No", "zF2", "zF3")])), c("Froude", "Froude^2", "Froude^3")))
POST_COR <- cor(as.data.frame(as_draws_df(FITS$M2))[, paste0("b_", c("z_LD_Ratio", "z_BD_Ratio", "z_LB_Ratio", "z_Froude_No", "zF2", "zF3"))])

# -----------------------------------------------------------------------------
# 11. POSTERIOR MANAGEMENT
# -----------------------------------------------------------------------------
hdi <- function(v, prob = .95) {
  v <- sort(v); k <- floor(prob * length(v)); w <- v[(k + 1):length(v)] - v[1:(length(v) - k)]
  i <- which.min(w); c(v[i], v[i + k])
}
kde_mode <- function(v) { d <- density(v); d$x[which.max(d$y)] }
manage <- function(v) {
  h <- hdi(v)
  c(mean = mean(v), median = median(v), mode_kde = kde_mode(v), sd = sd(v),
    eti_2.5 = unname(quantile(v, .025)), eti_97.5 = unname(quantile(v, .975)),
    hdi_low = h[1], hdi_high = h[2], P_gt_0 = mean(v > 0))
}
M0_RAW <- t(sapply(PRED, function(v) manage(d0[[paste0("b_z_", v)]] * SCALE$y[2] / SCALE[[v]][2])))
M0_RAW <- cbind(M0_RAW, OLS_raw_units = b_ols_raw[PRED])
BF <- d0$b_z_Froude_No * SCALE$y[2] / SCALE$Froude_No[2]

d6 <- as_draws_df(FITS$M6)
sl <- SCALE$ly[2]
EFF <- rbind(
  t(sapply(GEOM, function(v) { e <- 100 * (exp(d6[[paste0("b_z_", v)]] * sl) - 1); c(manage(e), P_abs_gt_1pct = mean(abs(e) > 1)) })),
  `Residual SD multiplier per +1 SD Froude` = { e <- exp(d6$b_sigma_z_Froude_No); c(manage(e), P_abs_gt_1pct = mean(e > 1.01)) },
  `Between-hull variation (%)` = { e <- 100 * (exp(d6$sd_hull__Intercept * sl) - 1); c(manage(e), P_abs_gt_1pct = mean(e > 1)) })

# kurva sigma(Froude) untuk hull tipikal
grid_df <- data.frame(Froude_No = seq(min(df$Froude_No), max(df$Froude_No), length.out = 60))
for (v in GEOM) grid_df[[v]] <- mean(df[[v]])
grid_df$Resistance <- 0; grid_df$hull_id <- 999
GRID <- FR$tf(grid_df)
SIG_GRID <- posterior_linpred(FITS$M6, newdata = GRID, dpar = "sigma", transform = TRUE, re_formula = NA)

# -----------------------------------------------------------------------------
# 12. PREDIKSI HULL BARU (posterior predictive, materi Bag. 24.6)
# -----------------------------------------------------------------------------
hull_geom <- df[!duplicated(df$hull_id), GEOM]
FROUDE_GRID <- sort(unique(df$Froude_No))
scen <- data.frame(Froude_No = FROUDE_GRID)
for (v in GEOM) scen[[v]] <- median(hull_geom[[v]])
scen$Resistance <- 0; scen$hull_id <- 999
SCEN <- FR$tf(scen)
predict_new_hull <- function(fit) {
  set.seed(SEED)
  to_raw(posterior_predict(fit, newdata = SCEN, allow_new_levels = TRUE, sample_new_levels = "gaussian"), "log", SCALE)
}
PRED_NEW <- predict_new_hull(FITS$M6)
MEAN_TYP <- to_raw(posterior_linpred(FITS$M6, newdata = SCEN, re_formula = NA), "log", SCALE)
PRED_TAB <- data.frame(Froude = FROUDE_GRID, pred_median = apply(PRED_NEW, 2, median),
                       PI_2.5 = apply(PRED_NEW, 2, quantile, .025), PI_97.5 = apply(PRED_NEW, 2, quantile, .975),
                       P_R_gt_5 = colMeans(PRED_NEW > 5), P_R_gt_20 = colMeans(PRED_NEW > 20))
I40 <- which(abs(FROUDE_GRID - 0.40) < 1e-9)

# -----------------------------------------------------------------------------
# 13. SENSITIVITAS: lab prior-data (analitik) & prior pada M6 (MCMC)
# -----------------------------------------------------------------------------
sigma2_known <- bn / (an - 1)
hull_order <- unique(mdl$hull_id)
SUBSETS <- list(`1 hull (n=14)` = hull_order[1], `3 hulls (n=42)` = hull_order[1:3], `22 hulls (n=308)` = hull_order)
PRIORS_LAB <- list(`Weak N(0, 1)` = c(0, 1), `Strong-correct N(0.81, 0.02)` = c(0.81, 0.02),
                   `Strong-wrong N(0.40, 0.02)` = c(0.40, 0.02))
LAB <- do.call(rbind, lapply(names(SUBSETS), function(sn) {
  idx <- mdl$hull_id %in% SUBSETS[[sn]]
  Xs <- X[idx, ]; ys <- y[idx]
  do.call(rbind, lapply(names(PRIORS_LAB), function(pn) {
    pr <- PRIORS_LAB[[pn]]
    m0v <- rep(0, p); m0v[p] <- pr[1]
    V0v <- diag(p); V0v[p, p] <- pr[2]^2
    Vni <- solve(V0v) + crossprod(Xs) / sigma2_known
    Vn_ <- solve(Vni)
    mn_ <- drop(Vn_ %*% (solve(V0v) %*% m0v + crossprod(Xs, ys) / sigma2_known))
    b_hat <- coef(lm(ys ~ Xs[, p]))[2]
    data.frame(data = sn, prior = pn, post_mean = mn_[p], post_sd = sqrt(Vn_[p, p]),
               data_only_slope = unname(b_hat), prior_pull = (1 / pr[2]^2) / Vni[p, p])
  }))
}))

SENS_FITS <- list(`Tighter (x0.5)` = fit_model("M6", s = 0.5, tag = "M6_prior_x0.5"),
                  `Base (x1)` = FITS$M6,
                  `Wider (x2.5)` = fit_model("M6", s = 2.5, tag = "M6_prior_x2.5"))
KEY <- c("b_z_LD_Ratio", "b_z_BD_Ratio", "b_z_LB_Ratio", "b_z_LCB", "b_z_Prismatic_Coef",
         "b_sigma_z_Froude_No", "sd_hull__Intercept")
SENS <- cache_rds("sensitivity", do.call(rbind, lapply(names(SENS_FITS), function(lab) {
  f <- SENS_FITS[[lab]]; dd <- as_draws_df(f)
  data.frame(prior = lab, param = KEY, median = sapply(KEY, function(k) median(dd[[k]])),
             lo = sapply(KEY, function(k) quantile(dd[[k]], .025)), hi = sapply(KEY, function(k) quantile(dd[[k]], .975)),
             P_R20_Fr040 = mean(predict_new_hull(f)[, I40] > 20),
             ELPD_LOO = loo_raw(f, "M6")$loo$estimates["elpd_loo", "Estimate"], row.names = NULL)
})))
NU <- as_draws_df(FITS$M5)$nu

# -----------------------------------------------------------------------------
# 14. AUDIT REPRODUCIBILITY: seed & urutan data
# -----------------------------------------------------------------------------
compare_runs <- function(fa, fb) {
  v <- MAIN_PAR(fa)
  sa <- summarise_draws(as_draws_array(fa, variable = v), "mean", "sd", "mcse_mean")
  sb <- summarise_draws(as_draws_array(fb, variable = v), "mean", "sd", "mcse_mean")
  data.frame(param = v, mean_a = sa$mean, mean_b = sb$mean, sd_a = sa$sd,
             z_mcse = (sb$mean - sa$mean) / sqrt(sa$mcse_mean^2 + sb$mcse_mean^2),
             pct_of_sd = 100 * abs(sb$mean - sa$mean) / sa$sd,
             expected_noise_pct = 100 * sqrt(2) * median(sa$mcse_mean / sa$sd))
}
SEEDS_EXTRA <- c(777, 999, 31, 123)
SEED_FITS <- c(list(`2026` = FITS$M6),
               lapply(setNames(SEEDS_EXTRA, SEEDS_EXTRA), function(s)
                 cache_rds(paste0("fit_M6_seed", s), update(FITS$M6, seed = s, refresh = 0, silent = 2, recompile = FALSE))))
FIT_SEED <- SEED_FITS[["777"]]
set.seed(SEED)
SHUFFLED <- mdl[sample(nrow(mdl)), ]
FIT_SHUF <- cache_rds("fit_M6_shuffled", update(FITS$M6, newdata = SHUFFLED, seed = SEED, refresh = 0,
                                                silent = 2, recompile = FALSE))
CMP_SEED <- compare_runs(FITS$M6, FIT_SEED)
CMP_ORDER <- compare_runs(FITS$M6, FIT_SHUF)
FLAGGED <- unique(c(CMP_SEED$param[abs(CMP_SEED$z_mcse) >= 3], CMP_ORDER$param[abs(CMP_ORDER$z_mcse) >= 3]))

# Audit lima seed: apakah MCSE yang dilaporkan sesuai dengan variasi nyata antar-run?
v_main <- MAIN_PAR(FITS$M6)
seed_sum <- lapply(SEED_FITS, function(f) summarise_draws(as_draws_array(f, variable = v_main), "mean", "mcse_mean"))
seed_means <- sapply(seed_sum, `[[`, "mean"); rownames(seed_means) <- v_main
MULTI_SEED <- data.frame(param = v_main, t(apply(seed_means, 1, range)),
                         between_run_sd = apply(seed_means, 1, sd),
                         mean_reported_mcse = rowMeans(sapply(seed_sum, `[[`, "mcse_mean")))
names(MULTI_SEED)[2:3] <- c("min_mean_5_seeds", "max_mean_5_seeds")
MULTI_SEED$ratio_sd_to_mcse <- MULTI_SEED$between_run_sd / MULTI_SEED$mean_reported_mcse
# Kesimpulan substantif antar seed: variasi antar-hull (%) dan probabilitas eksedans hull baru
SEED_CONCL <- do.call(rbind, lapply(names(SEED_FITS), function(s) {
  dd <- as_draws_df(SEED_FITS[[s]])
  data.frame(seed = s, between_hull_pct = median(100 * (exp(dd$sd_hull__Intercept * SCALE$ly[2]) - 1)),
             P_R20_Fr040 = mean(predict_new_hull(SEED_FITS[[s]])[, I40] > 20))
}))

# -----------------------------------------------------------------------------
# 15. CROSS-VALIDATION: per hull (grouped) vs acak per baris
# -----------------------------------------------------------------------------
logmeanexp <- function(x) { m <- max(x); m + log(mean(exp(x - m))) }
make_folds <- function(scheme) {
  if (scheme == "grouped") {
    hulls <- sort(unique(full$hull_id))
    fold_of_hull <- rep(1:5, length.out = length(hulls))            # 22 hull -> fold berisi 4-5 hull utuh
    fold_of_hull[match(full$hull_id, hulls)]
  } else {
    set.seed(SEED); sample(rep(1:5, length.out = nrow(full)))
  }
}
cv_eval <- function(name, scheme) cache_rds(paste0("cv_", name, "_", scheme), {
  sp <- SPECS[[name]]
  resp <- if (sp$resp == "raw") "y_z" else "ly_z"
  folds <- make_folds(scheme)
  recs <- list(); rhats <- c()
  for (k in 1:5) {
    fr <- make_frames(full[folds != k, ], full[folds == k, ])
    f <- update(FITS[[name]], newdata = fr$train, chains = CV_SAMPLER$chains, iter = CV_SAMPLER$iter,
                warmup = CV_SAMPLER$warmup, seed = SEED + k, refresh = 0, silent = 2, recompile = FALSE)
    rhats <- c(rhats, max(summarise_draws(as_draws_array(f, variable = MAIN_PAR(f)), "rhat")$rhat))
    ll <- log_lik(f, newdata = fr$test, allow_new_levels = TRUE, sample_new_levels = "gaussian")
    lpd <- apply(ll, 2, logmeanexp) + jacobian(fr$test, sp$resp, fr$sc)
    set.seed(SEED + k)
    yr <- to_raw(posterior_predict(f, newdata = fr$test, allow_new_levels = TRUE, sample_new_levels = "gaussian"),
                 sp$resp, fr$sc)
    recs[[k]] <- data.frame(fold = k, hull_id = fr$test$hull_id, Froude_No = fr$test$Froude_No,
                            y = fr$test$Resistance, pred_mean = colMeans(yr),
                            lo = apply(yr, 2, quantile, .025), hi = apply(yr, 2, quantile, .975), lpd = lpd)
  }
  r <- bind_rows(recs)
  list(pred = r, summary = data.frame(
    model = name, scheme = scheme, RMSE = sqrt(mean((r$y - r$pred_mean)^2)), MAE = mean(abs(r$y - r$pred_mean)),
    coverage_95 = mean(r$y >= r$lo & r$y <= r$hi), PI_width = mean(r$hi - r$lo),
    ELPD_per_obs = mean(r$lpd), ELPD_total = sum(r$lpd), max_rhat_folds = max(rhats)))
})
CV_PLAN <- list(c("M0", "grouped"), c("M1", "grouped"), c("M2", "grouped"), c("M3", "grouped"), c("M4", "grouped"),
                c("M6", "grouped"), c("M4", "row-random"), c("M6", "row-random"))
CV <- lapply(CV_PLAN, function(x) cv_eval(x[1], x[2]))
names(CV) <- sapply(CV_PLAN, paste, collapse = "_")
CV_TAB <- bind_rows(lapply(CV, `[[`, "summary"))

cv_key <- function(r) paste(r$hull_id, round(r$Froude_No, 3))
grouped_models <- CV_TAB$model[CV_TAB$scheme == "grouped"]
BEST_CV <- grouped_models[which.max(CV_TAB$ELPD_total[CV_TAB$scheme == "grouped"])]
ref <- CV[[paste0(BEST_CV, "_grouped")]]$pred
PAIRED <- bind_rows(lapply(grouped_models, function(m) {
  o <- CV[[paste0(m, "_grouped")]]$pred
  o <- o[match(cv_key(ref), cv_key(o)), ]
  d <- o$lpd - ref$lpd
  hs <- tapply(d, ref$hull_id, sum)
  data.frame(model = m, ELPD_total = sum(o$lpd), elpd_diff = sum(d),
             se_diff_obs = sd(d) * sqrt(length(d)), se_diff_hull = sd(hs) * sqrt(length(hs)))
})) %>% arrange(desc(ELPD_total)) %>% mutate(diff_over_se_hull = ifelse(se_diff_hull > 0, elpd_diff / se_diff_hull, NA))
OPTIMISM <- bind_rows(lapply(c("M4", "M6"), function(m) {
  g <- CV[[paste0(m, "_grouped")]]$pred; r <- CV[[paste0(m, "_row-random")]]$pred
  r <- r[match(cv_key(g), cv_key(r)), ]
  d <- r$lpd - g$lpd
  data.frame(model = m, ELPD_optimism = sum(d), se = sd(d) * sqrt(length(d)),
             RMSE_grouped = sqrt(mean((g$y - g$pred_mean)^2)), RMSE_row = sqrt(mean((r$y - r$pred_mean)^2)))
})) %>% mutate(optimism_over_se = ELPD_optimism / se)

# -----------------------------------------------------------------------------
# 16. CHECKLIST & MANIFEST REPRODUCIBILITY
# -----------------------------------------------------------------------------
CHECKS <- data.frame(Item = c(
  "Seed eksplisit tercatat", "Checksum SHA-256 data mentah tercatat", "Preprocessing scripted (make_frames)",
  "Standardisasi dihitung di dalam fold (tanpa kebocoran)", "Prior eksplisit + prior predictive check",
  "Komputasi divalidasi (analitik vs Gibbs vs HMC, < 3 MCSE)", "Likelihood manual = likelihood Stan",
  "Semua model: R-hat <= 1.01, ESS >= 400, 0 divergence", "Stabil terhadap seed (semua |z| < 3 MCSE)",
  "Stabil terhadap urutan data (semua |z| < 3 MCSE)",
  "Kesimpulan prediktif stabil antar 5 seed (rentang P(R>20 | Fr=0.40) < 0.05)",
  "Validasi per-hull dipakai untuk klaim prediktif", "Objek fit tersimpan (fits/*.rds)"),
  Pass = c(TRUE, nchar(RAW_SHA256) == 64, TRUE, TRUE, TRUE, max(abs(COMP$HMC_minus_Analytic_in_MCSE)) < 3,
           isTRUE(all.equal(LL_CHECK[["manual"]], LL_CHECK[["stan"]])), all(DIAG$verdict == "OK"),
           all(abs(CMP_SEED$z_mcse) < 3), all(abs(CMP_ORDER$z_mcse) < 3),
           diff(range(SEED_CONCL$P_R20_Fr040)) < 0.05, TRUE, file.exists("fits/M6.rds")))

# --------------------------------------------------------------------------
# Export CSV yang kompatibel dengan semua versi R:
# write.table()/write.csv() tidak dapat menulis kolom bertipe list.
# Hasil analisis tetap sama; bagian ini hanya mengubah cara file CSV disimpan.
# --------------------------------------------------------------------------
safe_csv <- function(x, file) {
  x <- as.data.frame(x, stringsAsFactors = FALSE)

  for (j in seq_along(x)) {
    if (is.factor(x[[j]])) {
      x[[j]] <- as.character(x[[j]])
    } else if (is.list(x[[j]])) {
      x[[j]] <- vapply(
        x[[j]],
        function(v) {
          if (length(v) == 0L || all(is.na(v))) return("")
          paste(as.character(unlist(v, use.names = FALSE)), collapse = "; ")
        },
        character(1)
      )
    }
  }

  utils::write.csv(x, file, row.names = FALSE, na = "")
}

safe_csv(CV_TAB, "cv_results.csv")
safe_csv(LOO_TAB, "model_comparison_loo.csv")
MANIFEST <- list(
  seed = SEED, raw_data_sha256 = RAW_SHA256, R = R.version.string,
  packages = sapply(c("brms", "rstan", "loo", "posterior", "dplyr", "splines"), function(x) as.character(packageVersion(x))),
  sampler = SAMPLER, cv_sampler = CV_SAMPLER, control = CONTROL,
  scaling_full_data = SCALE, diagnostics = DIAG, loo = LOO_TAB[, c("model", "elpd_loo", "elpd_diff", "se_diff")],
  cv = CV_TAB, multi_seed = MULTI_SEED, seed_conclusions = SEED_CONCL, checks = CHECKS)
jsonlite::write_json(MANIFEST, "reproducibility_manifest_R.json", pretty = TRUE, auto_unbox = TRUE, digits = NA)
\end{verbatim}

\end{document}
